\ifdefined\PRESENTATION\else\PassOptionsToClass{handout}{beamer}\fi\documentclass[aspectratio=169, lualatex]{beamer}
\makeatletter\def\input@path{{theme/}}\makeatother\usetheme{cipher}

\title{Applied Cryptography - 2.8: Zero-Knowledge Proofs}
\author{Nadim Kobeissi}
\subject{A comprehensive introduction to zero-knowledge proofs, from the Ali Baba cave and Schnorr identification through Sigma protocols, AND/OR proofs, the Fiat-Shamir transformation and its 2025 attack, circuits and zkVMs, to Zcash and zero-knowledge age verification.}
\keywords{zero-knowledge proofs, Schnorr protocol, Sigma protocols, cryptography, authentication, privacy, soundness, completeness, special soundness, simulation, Fiat-Shamir transformation, non-interactive proofs, OR proofs, AND proofs, discrete logarithm, forking lemma, circuits, zkVM, RISC Zero, Zcash, age verification, applied cryptography}
\institute{Symbolic Software}
\instituteimage{images/symbolicsoft-white.png}
\date{\today}
\coversubtitle{Summer 2026}
\coverpartname{Part 2: Real-World Cryptography}
\covertopicname{2.8: Zero-Knowledge Proofs}
\coverwebsite{https://appliedcryptography.page}

% Shared visual language for the mechanism diagrams (as in 2.7).
\colorlet{diagramgold}{cipheraccent!65!cipherdark}
\colorlet{diagramred}{red!65!cipherdark}
\tikzset{
	crypto diagram/.style={
		font=\sffamily\small,
		every node/.style={align=flush center, text=cipherdark, inner sep=4pt}
	},
	cd card/.style={draw=ciphergraydark, fill=cipherwhite, rounded corners=4pt, line width=0.7pt},
	cd green/.style={cd card, draw=cipherprimary!50, fill=cipherprimary!6!cipherwhite},
	cd gold/.style={cd card, draw=diagramgold!45, fill=cipheraccent!10!cipherwhite},
	cd red/.style={cd card, draw=diagramred!35, fill=diagramred!4!cipherwhite},
	cd heading/.style={font=\small\bfseries, text=cipherprimarydark},
	cd small/.style={font=\fontsize{8}{10}\selectfont},
	cd arrow/.style={-{Stealth[length=2.2mm, width=1.6mm]}, draw=cipherprimary, line width=1pt, rounded corners=3pt},
	cd block/.style={cd green, minimum width=1.65cm, minimum height=0.85cm, font=\small\bfseries},
	cd tag/.style={fill=cipherlight, inner sep=2pt, font=\fontsize{8}{10}\selectfont},
	% The cave: a ring of tunnel with an entrance at the bottom and a door at the top.
	cave rock/.style={draw=cipherdark!55, line width=13pt},
	cave floor/.style={draw=cipherlight, line width=10pt},
	cave person/.style={circle, minimum size=0.42cm, inner sep=0pt, font=\fontsize{8}{9}\selectfont\bfseries, text=cipherwhite}
}

% One cave panel: #1 = x offset, #2 = extra drawing (people, speech).
\newcommand{\cavepanel}[2]{
	\begin{scope}[xshift=#1 cm]
		\draw[cave rock] (0,0) circle (1.05);
		\draw[cave rock] (0,-1.05) -- (0,-2.05);
		\draw[cave floor] (0,0) circle (1.05);
		\draw[cave floor] (0,-1.0) -- (0,-2.1);
		\fill[diagramgold] (-0.09,0.87) rectangle (0.09,1.23);
		\node[cd small, text=diagramgold, font=\fontsize{7}{8}\selectfont\bfseries] at (0,1.5) {magic door};
		\node[font=\small\bfseries, text=cipherdark!55] at (-1.52,0) {A};
		\node[font=\small\bfseries, text=cipherdark!55] at (1.52,0) {B};
		#2
	\end{scope}
}

\begin{document}
\begin{frame}[plain]
	\titlepage
\end{frame}

\section{Introduction}

\begin{frame}{Zero-knowledge proofs: magic, but real}
	\begin{columns}[T]
		\begin{column}{0.48\textwidth}
			\textbf{The usual way to prove you know a password:}
			\begin{itemize}
				\item ``Here it is: \texttt{hunter2}''
				\item Now \textit{you} know the secret too
				\item You can impersonate me
				\item You can show anyone that I knew it
			\end{itemize}
		\end{column}
		\begin{column}{0.48\textwidth}
			\textbf{A zero-knowledge proof:}
			\begin{itemize}
				\item ``I won't tell you, but I'll convince you I know it''
				\item You end up certain that I know it
				\item You learn \textit{\textbf{nothing else}}: you can't impersonate me
				\item You can't even convince anyone else
			\end{itemize}
		\end{column}
	\end{columns}
	\vspace{0.4cm}
	\begin{alertblock}{The paradox}
		Convince you completely while teaching you nothing. Sounds impossible? Within the hour, you'll do it with pen and paper.
	\end{alertblock}
\end{frame}

\begin{frame}{The Ali Baba cave}
	\centering
	\begin{tikzpicture}[crypto diagram]
		\path[use as bounding box] (0,1.75) rectangle (13.5,-4.05);
		% Panel 1: Peggy walks in and picks a side; Victor waits outside.
		\cavepanel{2.1}{
			\draw[cipherprimary, dashed, line width=0.9pt, -{Stealth[length=1.8mm]}] (0,-1.95) -- (0,-1.05) arc (-90:-25:1.05);
			\node[cave person, fill=cipherprimary] at (-12:1.05) {P};
			\node[cave person, fill=diagramgold] at (0,-2.33) {V};
			\node[cd tag, text=cipherprimary] at (1.55,-1.05) {random\\side};
		}
		% Panel 2: Victor walks to the fork and shouts a side.
		\cavepanel{6.75}{
			\node[cave person, fill=cipherprimary] at (-12:1.05) {P};
			\node[cave person, fill=diagramgold] (v2) at (0,-1.5) {V};
			\node[cd card, font=\fontsize{8}{9}\selectfont\bfseries, inner sep=3pt] (shout) at (-1.55,-1.95) {``Come out\\of A!''};
			\draw[diagramgold, line width=0.7pt] (shout.east) -- (v2.west);
		}
		% Panel 3: Peggy comes out of the side Victor named, using the door.
		\cavepanel{11.4}{
			\draw[cipherprimary, line width=1.2pt, -{Stealth[length=2mm]}] (-12:1.05) arc (-12:200:1.05);
			\node[cave person, fill=cipherprimary] at (218:1.05) {P};
			\node[cave person, fill=diagramgold] at (0,-1.5) {V};
			\node[cd small, text=diagramgold, font=\fontsize{7}{8}\selectfont\bfseries] at (0,0.25) {says the\\magic word};
			\node[cd tag, text=cipherprimary, font=\fontsize{8}{9}\selectfont\bfseries] at (-1.5,-1.5) {A \mycheckmark};
		}
		\foreach \x/\step in {
				2.1/{\textbf{1.} Victor waits outside. Peggy takes path A or B at random.},
				6.75/{\textbf{2.} Victor walks to the fork and shouts a side.},
				11.4/{\textbf{3.} Peggy comes out that side, using the door if she must.}
			}
		\node[cd small, text width=4.1cm, font=\fontsize{8.5}{10.5}\selectfont, anchor=north] at (\x,-2.72) {\step};
		\node[cd small, text=cipherdark!65] at (6.75,-3.98) {Peggy (P) knows the magic word that opens the door. Victor (V) wants to be convinced without learning it.};
	\end{tikzpicture}
	\let\thefootnote\relax\footnote{Quisquater et al., ``How to Explain Zero-Knowledge Protocols to Your Children'', CRYPTO 1989.}
\end{frame}

\begin{frame}{The three properties, told by the cave}
	\begin{columns}[T]
		\begin{column}{0.32\textwidth}
			\begin{block}{1. Completeness}
				If the statement is true, an honest prover convinces an honest verifier.
			\end{block}
			\small\textbf{In the cave:} Peggy knows the word, so she can \textit{always} come out of the side Victor names.
		\end{column}
		\begin{column}{0.32\textwidth}
			\begin{block}{2. Soundness}
				If the statement is false, a cheating prover gets caught, except with negligible probability.
			\end{block}
			\small\textbf{In the cave:} without the word, Peggy must guess the shout in advance. Fooled with probability $1/2$ per round: one in a million after 20 rounds.
		\end{column}
		\begin{column}{0.32\textwidth}
			\begin{block}{3. Zero-knowledge}
				The verifier learns nothing beyond the fact that the statement is true.
			\end{block}
			\small\textbf{In the cave:} Victor only sees Peggy come out of the side he picked. He could have staged that scene himself\ldots
		\end{column}
	\end{columns}
\end{frame}

\begin{frame}{Why Victor's video convinces nobody else}
	\centering
	\begin{tikzpicture}[crypto diagram]
		\path[use as bounding box] (0,0) rectangle (13.5,-5.7);
		\foreach \y/\title/\subtitle/\lastx in {
				-0.8/{Victor's film}/{Peggy knows the word}/9.5,
				-2.3/{Mallory's raw footage}/{no word: she guesses}/12.65,
				-3.8/{Mallory's final cut}/{misses edited out}/9.5
			} {
				\node[anchor=west, font=\fontsize{9}{11}\selectfont\bfseries, text=cipherprimarydark, inner sep=0pt] at (0,\y+0.14) {\title};
				\node[anchor=west, cd small, text=cipherdark!65, inner sep=0pt] at (0,\y-0.24) {\subtitle};
				\fill[cipherdark!85, rounded corners=2pt] (3.1,\y+0.42) rectangle (\lastx,\y-0.42);
				\foreach \sx in {3.2,3.475,...,\lastx} {
						\fill[cipherlight] (\sx,\y+0.36) rectangle ++(0.11,-0.08);
						\fill[cipherlight] (\sx,\y-0.28) rectangle ++(0.11,-0.08);
					}
			}
		% Victor's genuine recording: every round a success.
		\foreach \i/\side in {0/A,1/B,2/B,3/A,4/B,5/A}
		\node[fill=cipherwhite, minimum width=0.9cm, minimum height=0.46cm, inner sep=0pt, font=\fontsize{8}{9}\selectfont\bfseries] at ({3.65+\i*1.05},-0.8) {\side\ \textcolor{cipherprimary}{\mycheckmark}};
		% Mallory's raw footage: she guessed wrong about half the time.
		\foreach \i/\side/\ok in {0/A/1,1/B/0,2/B/1,3/A/0,4/A/1,5/B/1,6/A/0,7/B/1,8/A/1} {
				\ifnum\ok=1
					\node[fill=cipherwhite, minimum width=0.9cm, minimum height=0.46cm, inner sep=0pt, font=\fontsize{8}{9}\selectfont\bfseries] at ({3.65+\i*1.05},-2.3) {\side\ \textcolor{cipherprimary}{\mycheckmark}};
				\else
					\node[fill=diagramred!12!cipherwhite, minimum width=0.9cm, minimum height=0.46cm, inner sep=0pt, font=\fontsize{8}{9}\selectfont\bfseries, text=diagramred] (miss\i) at ({3.65+\i*1.05},-2.3) {\side\ $\times$};
					\draw[diagramred, line width=0.8pt] (miss\i.north west) -- (miss\i.south east);
				\fi
			}
		% Mallory's final cut keeps only the successes.
		\foreach \i/\side in {0/A,1/B,2/A,3/B,4/B,5/A}
		\node[fill=cipherwhite, minimum width=0.9cm, minimum height=0.46cm, inner sep=0pt, font=\fontsize{8}{9}\selectfont\bfseries] at ({3.65+\i*1.05},-3.8) {\side\ \textcolor{cipherprimary}{\mycheckmark}};
		\draw[cd arrow, draw=diagramgold] (7.85,-2.75) -- (7.85,-3.35);
		\node[cd tag, text=diagramgold, anchor=west] at (8.0,-3.05) {cut every miss};
		\node[cd gold, cd small, text width=3.3cm] at (11.6,-0.8) {Carol: ``Nice film. How do I know it wasn't staged?''};
		\node[cd green, cd small, text width=3.3cm] at (11.6,-3.8) {Looks exactly like Victor's film};
		\node[cd small, text=cipherdark!65] at (6.75,-4.5) {Each frame: the side Victor shouted, and whether she came out of it.};
		\node[cd gold, text width=12.8cm, font=\fontsize{9}{11}\selectfont] at (6.75,-5.2) {A film that anyone can fake proves nothing, so watching it live taught Victor nothing he couldn't have staged himself. Mallory's editing room has a name: a \textbf{simulator}.};
	\end{tikzpicture}
\end{frame}

\begin{frame}{Today's route}
	\centering
	\begin{tikzpicture}[crypto diagram]
		\path[use as bounding box] (0,0) rectangle (13.5,-5.7);
		\node[anchor=west, cd heading] at (0,-0.3) {Start by proving you know one number. Finish by proving a whole program ran correctly.};
		\draw[ciphergraydark, line width=10pt, rounded corners=0.7cm] (0.2,-1.25) -- (13.1,-1.25) -- (13.1,-4.0) -- (0.2,-4.0);
		\draw[cipherwhite, line width=0.8pt, dash pattern=on 5pt off 4pt, rounded corners=0.7cm] (0.2,-1.25) -- (13.1,-1.25) -- (13.1,-4.0) -- (0.2,-4.0);
		\foreach \n/\x/\y/\name/\what in {
				1/1.45/-1.25/Schnorr/{Prove you know a private key, and nothing more},
				2/6.1/-1.25/Sigma protocols/{One pattern, many statements},
				3/10.75/-1.25/{AND, OR}/{Combine proofs; hide which branch is true},
				4/10.75/-4.0/Fiat-Shamir/{A hash plays the verifier: signatures, and a 2025 attack},
				5/6.1/-4.0/Circuits and zkVMs/{Prove any program ran correctly (Project D)},
				6/1.45/-4.0/Real world/{Zcash, and age checks on your phone}
			} {
				\node[circle, fill=cipherprimary, draw=cipherwhite, line width=1.2pt, minimum size=0.62cm, inner sep=0pt, font=\small\bfseries, text=cipherwhite] at (\x,\y) {\n};
				\node[font=\small\bfseries, text=cipherprimarydark] at (\x,\y-0.58) {\name};
				\node[cd small, text width=3.9cm, text=cipherdark!75, anchor=north] at (\x,\y-0.8) {\what};
			}
	\end{tikzpicture}
\end{frame}

\begin{frame}{From the cave to a private key}
	\centering
	\begin{tikzpicture}[crypto diagram]
		\path[use as bounding box] (0,0) rectangle (13.5,-5.7);
		\node[cd heading] at (3.2,-0.25) {In the cave};
		\node[cd heading] at (9.95,-0.25) {In Schnorr's protocol (1989)};
		\foreach \y/\cave/\schnorr in {
				-0.85/{The magic word}/{The private key $a$, behind the public key $A = g^a$},
				-1.5/{Peggy takes a random path}/{\textbf{Commit:} pick a random $y$, send $Y = g^y$},
				-2.15/{Victor shouts A or B}/{\textbf{Challenge:} a random $c \in \mathbb{Z}_n$},
				-2.8/{Peggy exits the named side}/{\textbf{Respond:} $r = y + ca \bmod n$},
				-3.45/{Victor sees her come out}/{\textbf{Check:} $g^r \overset{?}{=} Y \cdot A^c$},
				-4.1/{Fooled with prob.\ $1/2$ per round}/{Fooled with prob.\ $1/n \approx 2^{-256}$: one round is enough}
			} {
				\node[cd card, minimum width=5.4cm, minimum height=0.52cm, inner sep=2pt, font=\fontsize{9}{11}\selectfont] (cave) at (3.2,\y) {\cave};
				\node[cd green, minimum width=7.2cm, minimum height=0.52cm, inner sep=2pt, font=\fontsize{9}{11}\selectfont] (sch) at (9.95,\y) {\schnorr};
				\draw[cd arrow, draw=cipherprimary!60] (cave.east) -- (sch.west);
			}
		\node[cd gold, text width=12.6cm, font=\fontsize{9}{11}\selectfont] at (6.75,-5.05) {The same three moves: \textbf{commit}, \textbf{challenge}, \textbf{respond}. The cave has 2 possible challenges; Schnorr has $n$, a 256-bit prime. That's why one round suffices.};
	\end{tikzpicture}
\end{frame}

\begin{frame}{A toy group to play with}
	\centering
	\begin{tikzpicture}[crypto diagram]
		\node[anchor=east, font=\small\bfseries] at (1.4,0) {$k$};
		\node[anchor=east, font=\small\bfseries] at (1.4,-0.62) {$2^k \bmod 23$};
		\foreach \k/\v in {0/1,1/2,2/4,3/8,4/16,5/9,6/18,7/13,8/3,9/6,10/12} {
				\node[cd card, minimum width=0.98cm, minimum height=0.52cm, inner sep=0pt] at ({2+\k*1.05},0) {$\k$};
				\ifnum\k=7
					\node[cd gold, minimum width=0.98cm, minimum height=0.52cm, inner sep=0pt, font=\small\bfseries] at ({2+\k*1.05},-0.62) {$\v$};
				\else
					\node[cd green, minimum width=0.98cm, minimum height=0.52cm, inner sep=0pt, font=\small\bfseries] at ({2+\k*1.05},-0.62) {$\v$};
				\fi
			}
	\end{tikzpicture}
	\begin{columns}[T]
		\begin{column}{0.55\textwidth}
			\textbf{The rules} ($p = 23$, $g = 2$, $n = 11$):
			\begin{itemize}
				\item Elements live mod $23$; exponents live mod $11$, since $2^{11} \equiv 1$
				\item Multiplying elements adds exponents:\\$9 \cdot 12 = 2^{5+10} = 2^{15} = 2^{4} = 16$
				\item Inverses: $(2^k)^{-1} = 2^{11-k}$, so $13^{-1} = 2^{4} = 16$ (indeed $13 \cdot 16 = 208 \equiv 1$)
				\item Dividing exponents: $1/5 \equiv 9 \pmod{11}$, as $5 \cdot 9 = 45 \equiv 1$
			\end{itemize}
		\end{column}
		\begin{column}{0.42\textwidth}
			\begin{exampleblock}{Meet Alice}
				Private key $a = 7$. Public key $A = 2^7 = 13$.
			\end{exampleblock}
			\begin{alertblock}{Toy-sized on purpose}
				Here anyone can read $a$ off the table. Real groups have $n \approx 2^{256}$: no table, no shortcut (the discrete log problem, 1.7).
			\end{alertblock}
		\end{column}
	\end{columns}
\end{frame}

\begin{frame}{The Schnorr identification protocol}
	\vspace{-1.5cm}
	\begin{center}\resizebox{0.78\textwidth}{!}{
			\begin{msc}{}
				\setmscvalues{small}
				\drawframe{none}
				\setlength{\instdist}{8cm}
				\setlength{\instwidth}{2cm}
				\declinst{prover}{}{Prover}
				\declinst{verifier}{}{Verifier}
				\action*{$\begin{array}{l}
							\text{Knows generator}~g         \\
							\text{Knows prime}~n             \\
							\text{Knows secret}~a            \\
							y \twoheadleftarrow \mathbb{Z}_n \\
							Y = g^y
						\end{array}$}{prover}
				\action*{$\begin{array}{l}
							\text{Knows generator}~g                 \\
							\text{Knows prime}~n                     \\
							\text{Knows prover's public key}~A = g^a \\
						\end{array}$}{verifier}
				\nextlevel[7]
				\mess{$Y$}{prover}{verifier}
				\nextlevel[1]
				\action*{$c \twoheadleftarrow \mathbb{Z}_n$}{verifier}
				\nextlevel[2]
				\mess{$c$}{verifier}{prover}
				\nextlevel[1]
				\action*{$r = (y + ca) \bmod n$}{prover}
				\nextlevel[2]
				\mess{$r$}{prover}{verifier}
				\nextlevel[1]
				\action*{$g^r \overset{?}{\equiv} Y \cdot A^c$}{verifier}
				\nextlevel[2]
			\end{msc}
		}\end{center}
\end{frame}

\begin{frame}{Schnorr with toy numbers}
	\centering
	\begin{tikzpicture}[crypto diagram]
		\path[use as bounding box] (0,0) rectangle (13.5,-5.7);
		\node[cd green, minimum width=3.4cm, minimum height=1.1cm] at (1.9,-0.6) {\textbf{Prover (Alice)}\\\small knows $a = 7$};
		\node[cd card, minimum width=3.4cm, minimum height=1.1cm] at (11.6,-0.6) {\textbf{Verifier}\\\small knows $A = 13$};
		\draw[ciphergraydark, line width=1.2pt] (1.9,-1.15) -- (1.9,-4.2);
		\draw[ciphergraydark, line width=1.2pt] (11.6,-1.15) -- (11.6,-4.2);
		% Commit
		\draw[cd arrow] (1.9,-1.75) -- (11.6,-1.75);
		\node[cd tag, font=\small] at (6.75,-1.52) {$Y = 2^{y} = 2^{5} = 9$};
		\node[cd small, text=cipherdark!65, anchor=west] at (2.0,-2.0) {picks a random $y = 5$};
		% Challenge
		\draw[cd arrow, draw=diagramgold] (11.6,-2.7) -- (1.9,-2.7);
		\node[cd tag, font=\small] at (6.75,-2.47) {$c = 3$};
		\node[cd small, text=cipherdark!65, anchor=east] at (11.5,-2.95) {picks a random challenge};
		% Response
		\draw[cd arrow] (1.9,-3.65) -- (11.6,-3.65);
		\node[cd tag, font=\small] at (6.75,-3.42) {$r = y + ca = 5 + 3 \cdot 7 = 26 \equiv 4 \pmod{11}$};
		\node[cd small, text=cipherdark!65, anchor=west] at (2.0,-3.9) {uses the secret $a$};
		% Verification and completeness
		\node[cd green, minimum width=6.3cm, minimum height=1.05cm, font=\fontsize{9}{12}\selectfont] at (10.2,-4.95) {$g^r = 2^4 = 16$\\$Y \cdot A^c = 2^{5} \cdot 2^{7 \cdot 3} = 2^{26} = 2^{4} = 16$ \ \textcolor{cipherprimary}{\mycheckmark}};
		\node[cd card, minimum width=6.3cm, minimum height=1.05cm, font=\fontsize{9}{12}\selectfont] at (3.3,-4.95) {\textbf{Why it always works:}\\$g^{r} = g^{y + ca} = g^{y} \cdot (g^{a})^{c} = Y \cdot A^{c}$};
	\end{tikzpicture}
\end{frame}

\begin{frame}{Your turn: pass without knowing $a$}
	\begin{columns}[T]
		\begin{column}{0.48\textwidth}
			\textbf{You are Mallory.} You know $A = 13$ but not $a$. The verifier will pick $c$ uniformly from $\{0, \ldots, 10\}$. How often can you pass?
			\vspace{0.3cm}
			\pause

			\textbf{Attempt 1: commit honestly.}
			\begin{itemize}
				\item Send $Y = 2^y$ for a random $y$
				\item Then $c$ arrives, and you need $2^r = Y \cdot A^c = 2^{y + ca}$
				\item So $r = y + ca$: impossible without $a$ (in a real-sized group)
			\end{itemize}
		\end{column}
		\begin{column}{0.48\textwidth}
			\pause
			\textbf{Attempt 2: bet on the challenge.}
			\begin{itemize}
				\item Bet on $c^* = 2$: pick $r = 9$, send\\$Y = 2^r / A^{c^*} = 2^{9 - 14} = 2^{6} = 18$
				\item If $c = 2$ arrives: $2^9 = 6 = 18 \cdot 13^2$ \mycheckmark
				\item Any other $c$: stuck again
			\end{itemize}
			\begin{alertblock}{You win with probability $1/n$}
				Here $1/11$; with $n \approx 2^{256}$, never. But remember the trick: \textit{choose the challenge first, then solve for $Y$}.
			\end{alertblock}
		\end{column}
	\end{columns}
\end{frame}

\begin{frame}{A parable}
	\vfill
	{\rmfamily\itshape ``King Richard decides to hold an archery competition to identify the best archer in the land. He constructs a long wooden wall, with many targets painted on it. During the competition, he is stunned as Robin Hood stands 100 meters away from the wall and fires one arrow into the center of each of the targets. Of course, Robin Hood is declared the winner.''\par}
	\vspace{0.5cm}
	\pause
	{\rmfamily\itshape ``Later that day, the Sheriff of Nottingham arrives at the King's castle. The King describes Robin Hood's exceptional performance, pointing to the arrows still in the targets. The Sheriff is not impressed. ``Robin Hood is a fraud and a liar. Those arrows prove nothing,'' he says. He takes his own bow and fires a series of arrows into blank parts of the wall. He finds a bucket of paint and paints a target around each of his arrows. ``See,'' he says again, ``arrows in targets prove nothing!''\,''\par}
	\vfill
\end{frame}

\begin{frame}{The moral: targets first, or arrows first?}
	\centering
	\begin{tikzpicture}[crypto diagram,
			target/.pic={
					\fill[diagramred!80] (0,0) circle (0.42);
					\fill[cipherwhite] (0,0) circle (0.3);
					\fill[diagramred!80] (0,0) circle (0.17);
				},
			arrow in/.pic={
					\draw[cipherdark, line width=1.1pt] (0,0) -- (40:0.85);
					\draw[diagramgold, line width=1.4pt] (40:0.62) -- ++(0.02,0.16);
					\draw[diagramgold, line width=1.4pt] (40:0.62) -- ++(0.15,-0.05);
					\draw[diagramgold, line width=1.4pt] (40:0.74) -- ++(0.02,0.16);
					\draw[diagramgold, line width=1.4pt] (40:0.74) -- ++(0.15,-0.05);
				}
		]
		\path[use as bounding box] (0,0) rectangle (13.5,-5.7);
		\foreach \offset/\who/\one/\two in {
				0/{Robin Hood}/{Paint the targets}/{Shoot the arrows},
				7/{The Sheriff}/{Shoot the arrows}/{Paint the targets}
			} {
				\begin{scope}[xshift=\offset cm]
					\node[cd heading] at (3.25,-0.2) {\who};
					\fill[cipheraccent!25!cipherwhite, draw=diagramgold!50, rounded corners=3pt] (0.25,-0.55) rectangle (6.25,-2.55);
					\foreach \tx/\ty in {1.25/-1.65,3.25/-1.35,5.2/-1.75} {
							\pic at (\tx,\ty) {target};
							\pic at (\tx,\ty) {arrow in};
						}
					\node[cd card, cd small, anchor=west, minimum width=2.8cm] at (0.25,-3.0) {\textbf{1.} \one};
					\node[cd card, cd small, anchor=east, minimum width=2.8cm] at (6.25,-3.0) {\textbf{2.} \two};
				\end{scope}
			}
		\node[cd green, text width=5.7cm, font=\fontsize{9}{11}\selectfont] at (3.25,-3.95) {\textbf{Real prover:} commit to $Y$, \textit{then} the challenge $c$ arrives, then answer with $r$};
		\node[cd gold, text width=5.7cm, font=\fontsize{9}{11}\selectfont] at (10.25,-3.95) {\textbf{Mallory (and the simulator):} pick $c$ and $r$, \textit{then} solve for $Y$};
		\node[font=\fontsize{9}{11}\selectfont, text width=12.8cm] at (6.75,-5.1) {The two walls look identical. Only the \textbf{order of events} tells them apart, and a third party who sees the wall afterwards never saw the order.};
	\end{tikzpicture}
\end{frame}

\begin{frame}{The verifier could have written the transcript himself}
	\centering
	\begin{tikzpicture}[crypto diagram]
		\path[use as bounding box] (0,0) rectangle (13.5,-5.7);
		\node[anchor=west, font=\small\bfseries, text=cipherprimarydark] at (0,-0.78) {Real conversation};
		\node[anchor=west, cd small, text=cipherdark!65] at (0,-1.15) {with Alice, who knows $a$};
		\node[cd green, minimum width=2.4cm, minimum height=0.95cm] (r1) at (4.6,-0.95) {$Y = g^{y}$\\\scriptsize random $y$};
		\node[cd green, minimum width=2.4cm, minimum height=0.95cm] (r2) at (7.6,-0.95) {$c$\\\scriptsize random};
		\node[cd green, minimum width=2.4cm, minimum height=0.95cm] (r3) at (10.6,-0.95) {$r = y + ca$\\\scriptsize needs $a$};
		\node[anchor=west, font=\small\bfseries, text=cipherprimarydark] at (0,-2.38) {Victor, alone};
		\node[anchor=west, cd small, text=cipherdark!65] at (0,-2.75) {no secret at all};
		\node[cd gold, minimum width=2.4cm, minimum height=0.95cm] (s1) at (4.6,-2.55) {$c$\\\scriptsize random};
		\node[cd gold, minimum width=2.4cm, minimum height=0.95cm] (s2) at (7.6,-2.55) {$r$\\\scriptsize random};
		\node[cd gold, minimum width=2.4cm, minimum height=0.95cm] (s3) at (10.6,-2.55) {$Y = g^{r} / A^{c}$\\\scriptsize solved backwards};
		\draw[cd arrow] (r1) -- (r2);
		\draw[cd arrow] (r2) -- (r3);
		\draw[cd arrow, draw=diagramgold] (s1) -- (s2);
		\draw[cd arrow, draw=diagramgold] (s2) -- (s3);
		\node[font=\small\bfseries] at (12.75,-0.95) {$(Y, c, r)$};
		\node[font=\small\bfseries] at (12.75,-2.55) {$(Y, c, r)$};
		\node[cd small, text=cipherdark!65] at (7.6,-3.4) {time $\longrightarrow$};
		\node[cd card, text width=5.9cm, align=flush left, font=\fontsize{9}{11}\selectfont, anchor=north west] at (0,-3.8) {\textbf{Toy example:} Victor picks $c = 2$, $r = 9$, so $Y = 2^{9} / 13^{2} = 2^{9 - 14} = 2^{6} = 18$. The check passes: $2^9 = 6 = 18 \cdot 13^2$. (Mallory's bet, again!)};
		\node[cd green, text width=6.1cm, align=flush left, font=\fontsize{9}{11}\selectfont, anchor=north east] at (13.5,-3.8) {\textbf{Same distribution:} in both rows, $c$ and $r$ are uniform and independent, and $Y$ is the one value that passes the check. So a real transcript teaches Victor nothing he couldn't cook up alone.};
	\end{tikzpicture}
\end{frame}

\begin{frame}{Why $r$ leaks nothing: a one-time pad on a clock}
	\centering
	\begin{tikzpicture}[crypto diagram]
		\path[use as bounding box] (0,0) rectangle (13.5,-5.7);
		\node[anchor=west, font=\fontsize{9.5}{11}\selectfont] at (0,-0.25) {$r = y + ca \bmod n$ is the one-time pad from 1.2: key $y$, plaintext $ca$, addition mod $n$ for XOR.};
		% The clock of Z_11: start at ca, then spin by a uniformly random y.
		\begin{scope}[shift={(1.95,-2.45)}]
			\draw[ciphergraydark, line width=1.5pt] (0,0) circle (1.25);
			\foreach \k in {0,...,10} {
					\fill[cipherprimary] ({90-\k*360/11}:1.25) circle (2.2pt);
					\node[cd small, inner sep=0pt] at ({90-\k*360/11}:1.58) {\k};
				}
			\fill[diagramgold] ({90-10*360/11}:1.25) circle (4pt);
			\draw[cd arrow, draw=diagramgold, line width=1.2pt] ({90-10*360/11}:0.95) arc ({90-10*360/11}:{90-10*360/11-150}:0.95);
			\node[cd small, text=diagramgold, font=\fontsize{8}{9}\selectfont\bfseries] at (0,0.1) {start at $ca$};
			\node[cd small, text=diagramgold, font=\fontsize{8}{9}\selectfont\bfseries] at (0,-0.35) {spin by $+y$};
		\end{scope}
		% Two different secrets, same distribution of r.
		\node[anchor=east, cd small, font=\fontsize{8.5}{10}\selectfont\bfseries] at (6.05,-1.3) {$y$};
		\node[anchor=east, cd small, font=\fontsize{8.5}{10}\selectfont\bfseries] at (6.05,-1.92) {$r$ if $a = 7$};
		\node[anchor=east, cd small, font=\fontsize{8.5}{10}\selectfont\bfseries] at (6.05,-2.54) {$r$ if $a = 2$};
		\foreach \k in {0,...,10} {
				\pgfmathtruncatemacro{\ra}{mod(\k+10,11)}
				\pgfmathtruncatemacro{\rb}{mod(\k+6,11)}
				\node[cd card, minimum width=0.56cm, minimum height=0.5cm, inner sep=0pt, font=\small] at ({6.4+\k*0.6},-1.3) {\k};
				\node[cd green, minimum width=0.56cm, minimum height=0.5cm, inner sep=0pt, font=\small] at ({6.4+\k*0.6},-1.92) {\ra};
				\node[cd gold, minimum width=0.56cm, minimum height=0.5cm, inner sep=0pt, font=\small] at ({6.4+\k*0.6},-2.54) {\rb};
			}
		\node[cd small, text width=7.2cm, align=flush left, anchor=north west] at (6.1,-2.95) {With $c = 3$: $ca \equiv 10$ if $a = 7$, and $ca \equiv 6$ if $a = 2$. Each row is a reshuffle of $0, \ldots, 10$: every $r$ appears exactly once. \textbf{If $y$ is uniform, $r$ is uniform, whatever $a$ is.}};
		\node[cd gold, text width=12.8cm, font=\fontsize{9}{11}\selectfont] at (6.75,-5.05) {\textbf{And $Y = g^y$?} Once $c$ and $r$ are fixed, $Y$ is pinned down: it's the only value with $g^r = Y \cdot A^c$, and anyone can compute it as $g^r / A^c$. So $Y$ adds nothing either.};
	\end{tikzpicture}
\end{frame}

\begin{frame}{In other words\ldots}
	\begin{center}
		\sslinked{
			\sslibrary{}{schnorr-real}{
				\sslibrarysubroutine{schnorr.transcript}{a}{
					$A := g^a$ \\
					$y \twoheadleftarrow \mathbb{Z}_n$ \\
					$Y := g^y$ \\
					$c \twoheadleftarrow \mathbb{Z}_n$ \\
					$r := (y + ca) \bmod n$ \\
					return $(A, Y, c, r)$
				}{1}
			}{1}
		}{\interchangeable{}}{
			\sslibrary{}{schnorr-fake}{
				\sslibrarysubroutine{schnorr.transcript}{a}{
					$A := g^a$ \\
					$c \twoheadleftarrow \mathbb{Z}_n$ \\
					$r \twoheadleftarrow \mathbb{Z}_n$ \\
					$Y := g^r \cdot (A^c)^{-1}$ \\
					return $(A, Y, c, r)$
				}{1}
			}{1}
		}
	\end{center}
\end{frame}

\begin{frame}{Proof of indistinguishability}{Step 1}
	\begin{columns}[c]
		\begin{column}{0.5\textwidth}
			\begin{itemize}[<+->]
				\item $c$ doesn't depend on $y$, so we can reorder the lines and choose $c$ first.
				\item Nothing about the output changes: this is pure bookkeeping.
			\end{itemize}
		\end{column}
		\begin{column}{0.5\textwidth}
			\begin{center}
				\sslibrary{}{schnorr-real}{
					\sslibrarysubroutine{schnorr.transcript}{a}{
						$A := g^a$ \\
						$c \twoheadleftarrow \mathbb{Z}_n$ \\
						\hl{$y \twoheadleftarrow \mathbb{Z}_n$} \\
						$r := (y + ca) \bmod n$ \\
						\hl{$Y := g^y$} \\
						return $(A, Y, c, r)$
					}{1}
				}{1}
			\end{center}
		\end{column}
	\end{columns}
\end{frame}

\begin{frame}{Proof of indistinguishability}{Step 2}
	\begin{columns}[c]
		\begin{column}{0.5\textwidth}
			\begin{itemize}[<+->]
				\item $r$ is a one-time pad encryption of $ca$ under key $y$: the clock from before.
				\item So we get the same distribution by choosing $r$ uniformly and solving for $y$.
			\end{itemize}
		\end{column}
		\begin{column}{0.5\textwidth}
			\begin{center}
				\sslibrary{}{schnorr-real}{
					\sslibrarysubroutine{schnorr.transcript}{a}{
						$A := g^a$ \\
						$c \twoheadleftarrow \mathbb{Z}_n$ \\
						\hl{$r \twoheadleftarrow \mathbb{Z}_n$} \\
						\hl{$y := (r - ca) \bmod n$} \\
						$Y := g^y$ \\
						return $(A, Y, c, r)$
					}{1}
				}{1}
			\end{center}
		\end{column}
	\end{columns}
\end{frame}

\begin{frame}{Proof of indistinguishability}{Step 3}
	\begin{columns}[c]
		\begin{column}{0.5\textwidth}
			\begin{itemize}[<+->]
				\item Rewrite $Y$ without $y$:
				      \begin{align*}
					      Y & = g^{y} = g^{r - ca} = g^{r} \cdot (g^{ca})^{-1} \\
					        & = g^{r} \cdot (A^{c})^{-1}
				      \end{align*}
				      (toy check: $2^{9 - 14} = 2^{9} \cdot (2^{14})^{-1} = 6 \cdot 8^{-1} = 6 \cdot 3 = 18$)
				\item Now $y$ is never used. Delete its line, and this is exactly $\lib{}{schnorr-fake}$. $\blacksquare$
			\end{itemize}
		\end{column}
		\begin{column}{0.5\textwidth}
			\begin{center}
				\sslibrary{}{schnorr-fake}{
					\sslibrarysubroutine{schnorr.transcript}{a}{
						$A := g^a$ \\
						$c \twoheadleftarrow \mathbb{Z}_n$ \\
						$r \twoheadleftarrow \mathbb{Z}_n$ \\
						$y := (r - ca) \bmod n$ \\
						\hl{$Y := g^r \cdot (A^c)^{-1}$} \\
						return $(A, Y, c, r)$
					}{1}
				}{1}
			\end{center}
		\end{column}
	\end{columns}
\end{frame}

\begin{frame}{Another parable}
	\vfill
	{\rmfamily\itshape ``Alice, whose public key is $A = g^a$, contacts Bob and uses Schnorr's protocol to prove that she knows the secret key $a$. Bob is convinced that he is talking to Alice, and he writes down everything he sees: $(A, Y, c, r)$.''\par}
	\vspace{0.5cm}
	\pause
	{\rmfamily\itshape ``Later that day, Charlie arrives at Bob's house. Bob mentions that he talked to Alice earlier in the day, and he shows Charlie the accepting transcript $(A,Y,c,r)$. ``That proves nothing,'' says Charlie, and he proceeds to generate his own accepting transcript $(A,Y',c',r')$, even though he doesn't know $a$.''\par}
	\vspace{0.5cm}
	\pause
	\begin{alertblock}{The moral: Schnorr is deniable}
		Bob is convinced, but his transcript is worthless as evidence. Hold on to this: Fiat-Shamir will take it away.
	\end{alertblock}
	\vfill
\end{frame}

\section{Sigma Protocols}

\begin{frame}{Zero-knowledge is not just for logging in}
	\begin{columns}[c]
		\begin{column}{0.55\textwidth}
			\textbf{Statements we'd like to prove in zero knowledge:}
			\begin{itemize}
				\item \textit{``I know a solution to this Sudoku''}
				\item \textit{``This encrypted number is between 0 and 100''} (a range proof)
				\item \textit{``This encrypted vote is a yes or a no''}
				\item \textit{``I ran this program, and this is its output''}
				\item \textit{``My answer to your Battleship shot is honest''}
			\end{itemize}
			\vspace{0.2cm}
			Schnorr handles one statement: ``I know $a$ with $g^a = A$''. Time to generalize.
		\end{column}
		\begin{column}{0.4\textwidth}
			\imagewithcaption[0.8\textwidth]{battleship.jpg}{Project D: the server proves every hit or miss is honest, without revealing where its ships are. Source: Hasbro}
		\end{column}
	\end{columns}
\end{frame}

\subsection{Sigma Protocols: High-Level Overview}

\begin{frame}{From Schnorr to Sigma, a general framework}
	\begin{columns}[c]
		\begin{column}{0.5\textwidth}
			\textbf{Schnorr was just the beginning!}
			\begin{itemize}
				\item Schnorr proves: ``I know $a$ such that $g^a = A$''
				\item \textbf{Sigma protocols}: the same moves, for other statements
			\end{itemize}
			\textbf{The pattern:}
			\begin{itemize}
				\item Public \textbf{instance}: $X$ (everyone sees it)
				\item Private \textbf{witness}: $W$ (makes it true)
				\item \textbf{Relation}: $R(X, W) = \text{true}$
			\end{itemize}
		\end{column}
		\begin{column}{0.5\textwidth}
			\begin{exampleblock}{Examples of relations}
				\begin{itemize}
					\item ``I know $a$ such that $g^a = X$''
					      \begin{itemize}
						      \item Witness: $a$
					      \end{itemize}
					\item ``I know $(a, b)$ such that $g^a h^b = X$''
					      \begin{itemize}
						      \item Witness: $(a, b)$
					      \end{itemize}
					\item ``I know a factorization of $X$''
					      \begin{itemize}
						      \item Witness: $(p, q)$ where $X = p \cdot q$
					      \end{itemize}
					\item ``I know a 3-coloring of the graph $X$''
					      \begin{itemize}
						      \item Witness: the coloring
					      \end{itemize}
				\end{itemize}
			\end{exampleblock}
		\end{column}
	\end{columns}
\end{frame}

\begin{frame}{Sigma protocol structure}
	\vspace{-1.5cm}
	\begin{center}\resizebox{0.78\textwidth}{!}{
			\begin{msc}{}
				\setmscvalues{small}
				\drawframe{none}
				\setlength{\instdist}{8cm}
				\setlength{\instwidth}{2cm}
				\declinst{prover}{}{Prover}
				\declinst{verifier}{}{Verifier}
				\action*{$\begin{array}{l}
							\text{Knows instance}~X                \\
							\text{Knows witness}~W                 \\
							\text{Such that}~R(X, W) = \text{true} \\
						\end{array}$}{prover}
				\action*{$\begin{array}{l}
							\text{Knows instance}~X \\
						\end{array}$}{verifier}
				\nextlevel[4]
				\action*{$(Y, \sigma) \leftarrow \text{Commit}(X, W)$}{prover}
				\nextlevel[2]
				\mess{$Y$ (commitment)}{prover}{verifier}
				\nextlevel[1]
				\action*{$c \twoheadleftarrow \mathcal{C}$}{verifier}
				\nextlevel[2]
				\mess{$c$ (challenge)}{verifier}{prover}
				\nextlevel[1]
				\action*{$r \leftarrow \text{Respond}(\sigma, c)$}{prover}
				\nextlevel[2]
				\mess{$r$ (response)}{prover}{verifier}
				\nextlevel[1]
				\action*{$\text{Check}(X, Y, c, r) \overset{?}{=} \text{true}$}{verifier}
				\nextlevel[2]
			\end{msc}
		}\end{center}
\end{frame}

\begin{frame}{Why ``Sigma''? The shape of the protocol}
	\begin{columns}[c]
		\begin{column}{0.5\textwidth}
			\centering
			\begin{tikzpicture}[crypto diagram]
				\node[cd green, minimum width=1.7cm] at (0,0.45) {\textbf{Prover}};
				\node[cd card, minimum width=1.7cm] at (3.6,0.45) {\textbf{Verifier}};
				\draw[ciphergraydark, line width=1.2pt] (0,0) -- (0,-4.1);
				\draw[ciphergraydark, line width=1.2pt] (3.6,0) -- (3.6,-4.1);
				\draw[cd arrow, line width=1.6pt] (0,-0.3) -- node[cd tag, font=\small\bfseries, sloped, above=1pt] {$Y$} (3.6,-1.5);
				\draw[cd arrow, line width=1.6pt, draw=diagramgold] (3.6,-1.5) -- node[cd tag, font=\small\bfseries, sloped, above=1pt] {$c$} (0,-2.7);
				\draw[cd arrow, line width=1.6pt] (0,-2.7) -- node[cd tag, font=\small\bfseries, sloped, above=1pt] {$r$} (3.6,-3.9);
				\node[font=\small, text=cipherdark!55] at (4.55,-2.1) {$\approx$};
				\node[font=\fontsize{64}{64}\selectfont, text=cipherprimary!45] at (5.35,-2.05) {$\Sigma$};
			\end{tikzpicture}
		\end{column}
		\begin{column}{0.5\textwidth}
			\textbf{The Greek letter $\Sigma$ (Sigma)}
			\begin{itemize}
				\item Three messages, alternating direction
				\item Drawn over time, they zigzag across the page
				\item Squint, and the zigzag is a $\Sigma$
			\end{itemize}
			\begin{alertblock}{Always this pattern}
				\begin{enumerate}
					\item Prover $\rightarrow$ Verifier: \textbf{commitment}
					\item Verifier $\rightarrow$ Prover: \textbf{challenge}
					\item Prover $\rightarrow$ Verifier: \textbf{response}
				\end{enumerate}
			\end{alertblock}
		\end{column}
	\end{columns}
\end{frame}

\begin{frame}{Formal definition of Sigma protocols}
	\begin{block}{Definition: Sigma Protocol}
		A Sigma protocol consists of the following components:
		\begin{itemize}
			\item \textbf{Commit}: A randomized algorithm that takes instance $X$ and witness $W$ as input, outputs commitment $Y$ and state $\sigma$
			\item $\mathcal{C}$: A set of verifier challenges
			\item \textbf{Respond}: A deterministic algorithm that takes state $\sigma$ and challenge $c \in \mathcal{C}$ as input, outputs response $r$
			\item \textbf{Check}: A deterministic algorithm that takes transcript $(X, Y, c, r)$ as input, outputs true/false
		\end{itemize}
	\end{block}
	\textbf{Note:} the state $\sigma$ is the prover's private memory between moves. In Schnorr, it's the mask $y$.
\end{frame}

\begin{frame}{Example: Schnorr as a Sigma protocol}
	\begin{columns}[c]
		\begin{column}{0.5\textwidth}
			\textbf{Instance and witness:}
			\begin{itemize}
				\item Instance: $X = A = g^a$
				\item Witness: $W = a$
				\item Relation: $g^W = X$
			\end{itemize}
			\textbf{The algorithms:}
			\begin{itemize}
				\item \textbf{Commit}$(A, a)$:
				      \begin{itemize}
					      \item $y \twoheadleftarrow \mathbb{Z}_n$
					      \item $Y = g^y$
					      \item Return $(Y, \sigma = y)$
				      \end{itemize}
			\end{itemize}
		\end{column}
		\begin{column}{0.5\textwidth}
			\begin{itemize}
				\item $\mathcal{C} = \mathbb{Z}_n$ (all possible challenges)
				\item \textbf{Respond}$(\sigma = y, c)$:
				      \begin{itemize}
					      \item Return $r = (y + ca) \bmod n$
				      \end{itemize}
				\item \textbf{Check}$(A, Y, c, r)$:
				      \begin{itemize}
					      \item Return $(g^r = Y \cdot A^c)$
				      \end{itemize}
			\end{itemize}
			\begin{exampleblock}{Fits the framework!}
				Schnorr is just one example of a Sigma protocol
			\end{exampleblock}
		\end{column}
	\end{columns}
\end{frame}

\begin{frame}{Example: proving knowledge of a representation}
	\begin{columns}[c]
		\begin{column}{0.5\textwidth}
			\textbf{The problem:}
			\begin{itemize}
				\item Given: $X = g^a h^b$ (two generators)
				\item Prove: ``I know $(a, b)$''
				\item Without revealing $(a, b)$!
			\end{itemize}
			\textbf{Why is this useful?}
			\begin{itemize}
				\item A Pedersen commitment $X = g^m h^s$ hides a value $m$: this proves you can open it
				\item Anonymous credentials, mix networks, electronic voting
			\end{itemize}
		\end{column}
		\begin{column}{0.5\textwidth}
			\textbf{The Sigma protocol (Okamoto, 1992):}
			\begin{itemize}
				\item \textbf{Commit}: pick random $y, z$; send $Y = g^y h^z$ \\(state: $\sigma = (y, z)$)
				\item \textbf{Respond}: $u = y + ca$ and $v = z + cb$
				\item \textbf{Check}: $g^u h^v \overset{?}{=} Y \cdot X^c$
			\end{itemize}
			\begin{exampleblock}{Schnorr, twice over}
				One mask per secret, one shared challenge. Completeness: $g^{y + ca} h^{z + cb} = g^y h^z \cdot (g^a h^b)^c$.
			\end{exampleblock}
		\end{column}
	\end{columns}
\end{frame}

\begin{frame}{How far do Sigma protocols reach?}
	\begin{columns}[c]
		\begin{column}{0.5\textwidth}
			\textbf{Statements with algebraic structure (efficient):}
			\begin{itemize}
				\item ``I know $a$ with $g^a = A$'' (Schnorr)
				\item ``I know $(a, b)$ with $g^a h^b = X$'' (representation)
				\item ``These two ciphertexts encrypt the same value'' (an AND proof, coming up)
				\item ``This ciphertext encrypts 0 or 1'' (an OR proof, coming up)
				\item ``I shuffled these ciphertexts correctly'' (mix networks)
			\end{itemize}
		\end{column}
		\begin{column}{0.5\textwidth}
			\textbf{Statements without it (possible, but painful):}
			\begin{itemize}
				\item ``I know a 3-coloring of this graph''
				\item ``I know a Hamiltonian cycle in this graph''
			\end{itemize}
			\begin{alertblock}{Universal, in principle}
				Hamiltonian cycle is NP-complete and has a Sigma protocol (Blum, 1986), so \textit{every} NP statement has one, via a reduction. In practice that route is hopelessly slow. Circuits, later today, are the fix.
			\end{alertblock}
		\end{column}
	\end{columns}
\end{frame}

\subsection{Formal Properties of Sigma Protocols}

\begin{frame}{The three properties, now made precise}
	\begin{columns}[c]
		\begin{column}{0.5\textwidth}
			\textbf{Reminder} from the cave:
			\begin{enumerate}
				\item \textbf{Completeness:} honest provers always succeed
				\item \textbf{Soundness:} dishonest provers almost always fail
				\item \textbf{Zero-knowledge:} nothing leaks about the witness
			\end{enumerate}
		\end{column}
		\begin{column}{0.5\textwidth}
			\textbf{What's new for Sigma protocols:}
			\begin{itemize}
				\item Soundness comes in a stronger, very concrete form: \textbf{special soundness}, with an explicit extractor
				\item Zero-knowledge comes in a weaker form: \textbf{honest-verifier} zero-knowledge
				\item We'll define each one, then check it on Schnorr
			\end{itemize}
		\end{column}
	\end{columns}
\end{frame}

\begin{frame}{Completeness: honest provers always win}
	\begin{columns}[c]
		\begin{column}{0.5\textwidth}
			\begin{block}{Definition: Completeness}
				For any instance $X$ and witness $W$ with $R(X,W) = \text{true}$, if the prover follows the protocol honestly, the verifier \textit{always} accepts.
			\end{block}
			\textbf{In other words:}
			\begin{itemize}
				\item If you actually know the secret, you can always prove it
				\item No ``false negatives'': legitimate users never get locked out
			\end{itemize}
		\end{column}
		\begin{column}{0.5\textwidth}
			\textbf{Schnorr, for any $y$ and any $c$:}
			\begin{align*}
				g^r & = g^{y + ca}                   \\
				    & = g^y \cdot g^{ca}             \\
				    & = Y \cdot (g^a)^c              \\
				    & = Y \cdot A^c \quad \mycheckmark
			\end{align*}
			\begin{exampleblock}{Perfect completeness}
				Honest provers succeed with probability 1. (Toy run: $16 = 16$.)
			\end{exampleblock}
		\end{column}
	\end{columns}
\end{frame}

\begin{frame}{Soundness isn't enough: knowledge soundness}
	\begin{columns}[c]
		\begin{column}{0.5\textwidth}
			\textbf{Plain soundness:}
			\begin{itemize}
				\item ``If the statement is false, no prover can convince the verifier''
				\item But Schnorr's statement is ``there exists $a$ with $g^a = A$''
				\item In a cyclic group, \textit{every} $A$ has such an $a$ (look at the toy table: every element appears)
				\item Always true, so plain soundness guarantees nothing here!
			\end{itemize}
		\end{column}
		\begin{column}{0.5\textwidth}
			\textbf{Knowledge soundness:}
			\begin{itemize}
				\item ``If a prover convinces the verifier, a witness\footnote{A witness is the secret information that makes a statement true: the private key behind a public key, or the solution to a puzzle.} can be \textit{extracted} from that prover''
				\item The prover doesn't just show that $a$ exists: it must \textit{know} $a$
				\item The guarantee authentication and signatures actually need
				\item Next: the concrete version, \textbf{special soundness}
			\end{itemize}
		\end{column}
	\end{columns}
\end{frame}

\begin{frame}{One answer hides $a$; two answers reveal it}
	\centering
	\begin{tikzpicture}[crypto diagram]
		\path[use as bounding box] (0,0) rectangle (13.5,-5.7);
		\node[anchor=west, font=\fontsize{9.5}{11}\selectfont] at (0,-0.25) {Read the response as a line, $r = y + a \cdot c$: the \textbf{intercept} $y$ is the mask, the \textbf{slope} $a$ is the secret.};
		% Left: one transcript.
		\node[cd heading] at (3.35,-0.85) {One transcript $(Y, c, r)$};
		\begin{scope}[shift={(0.8,-4.35)}]
			\draw[-{Stealth[length=2mm]}, cipherdark!60, line width=0.8pt] (0,0) -- (5.2,0) node[right, inner sep=1pt] {$c$};
			\draw[-{Stealth[length=2mm]}, cipherdark!60, line width=0.8pt] (0,0) -- (0,3.05) node[above, inner sep=1pt] {$r$};
			\begin{scope}
				\clip (0,0) rectangle (5,2.95);
				\foreach \s in {-0.3,0.15,0.45,0.7}
				\draw[diagramgold, dashed, line width=0.9pt] (0,{1.8-2.4*\s}) -- (5,{1.8+2.6*\s});
			\end{scope}
			\foreach \s in {-0.3,0.15,0.45,0.7}
			\fill[diagramgold] (0,{1.8-2.4*\s}) circle (2pt);
			\node[cd small, text=diagramgold, anchor=east] at (-0.05,1.45) {$y'$?};
			\fill[cipherprimary] (2.4,1.8) circle (3pt);
			\node[cd small, font=\fontsize{8}{9}\selectfont\bfseries, anchor=south east] at (2.35,1.85) {$(c, r)$};
		\end{scope}
		\node[cd small, text width=5.6cm, anchor=north] at (3.35,-4.6) {Every slope $a'$ fits, with its own intercept $y' = r - a'c$. The verifier never sees $y$ (only $Y = g^y$), so no $a'$ is ruled out.};
		% Right: two transcripts with the same Y.
		\node[cd heading] at (10.15,-0.85) {Two transcripts, same $Y$};
		\begin{scope}[shift={(7.6,-4.35)}]
			\draw[-{Stealth[length=2mm]}, cipherdark!60, line width=0.8pt] (0,0) -- (5.2,0) node[right, inner sep=1pt] {$c$};
			\draw[-{Stealth[length=2mm]}, cipherdark!60, line width=0.8pt] (0,0) -- (0,3.05) node[above, inner sep=1pt] {$r$};
			\draw[cipherprimary, line width=1.3pt] (0,0.75) -- (5,2.95);
			\fill[cipherprimary] (0,0.75) circle (2.5pt);
			\node[cd small, text=cipherprimary, anchor=east] at (-0.05,0.75) {$y$};
			\draw[cipherdark!50, dashed] (1.2,1.278) -- (4.1,1.278) -- (4.1,2.554);
			\node[cd small, anchor=north] at (2.65,1.278) {$c - c'$};
			\node[cd small, anchor=west] at (4.12,1.9) {$r - r'$};
			\fill[cipherprimary] (1.2,1.278) circle (3pt);
			\fill[cipherprimary] (4.1,2.554) circle (3pt);
			\node[cd small, font=\fontsize{8}{9}\selectfont\bfseries, anchor=south east] at (1.15,1.33) {$(c', r')$};
			\node[cd small, font=\fontsize{8}{9}\selectfont\bfseries, anchor=south east] at (4.05,2.6) {$(c, r)$};
		\end{scope}
		\node[cd small, text width=5.6cm, anchor=north] at (10.15,-4.6) {Same $Y$ means same intercept. Two points pin down the line, and its slope is the secret: $a = (r - r')/(c - c')$.};
	\end{tikzpicture}
\end{frame}

\begin{frame}{Special soundness: formal definition}
	\begin{block}{Definition: Special Soundness}
		A Sigma protocol has \textbf{special soundness} if there exists an efficient algorithm \textbf{Extract} such that:
		\begin{itemize}
			\item Given two accepting transcripts: $(X, Y, c, r)$ and $(X, Y, c', r')$
			\item With the same instance $X$ and commitment $Y$
			\item But different challenges $c \neq c'$
			\item \textbf{Extract}$(X, Y, c, r, c', r')$ outputs witness $W$
			\item Such that $R(X, W) = \text{true}$
		\end{itemize}
	\end{block}
	\textbf{Interpretation:} being ready for \textit{two} different challenges on the same commitment means you know the secret. Why ``special''? It's soundness with a concrete recipe for the extractor.
\end{frame}

\begin{frame}{Example: extracting from Schnorr}
	\begin{columns}[c]
		\begin{column}{0.5\textwidth}
			\textbf{Two accepting transcripts, same $Y$:}
			\begin{itemize}
				\item $(Y, c, r)$, valid since $g^r = Y \cdot A^c$
				\item $(Y, c', r')$, valid since $g^{r'} = Y \cdot A^{c'}$
				\item Different challenges: $c \neq c'$, so $(c - c')^{-1} \bmod n$ exists ($n$ is prime)
			\end{itemize}
			\begin{exampleblock}{Toy check (Alice, $Y = 9$)}
				$(c, r) = (3, 4)$ and $(c', r') = (8, 6)$:\\
				$a = \frac{4 - 6}{3 - 8} = \frac{-2}{-5} = 2 \cdot 5^{-1} = 2 \cdot 9 = 18 \equiv 7$ \mycheckmark
			\end{exampleblock}
		\end{column}
		\begin{column}{0.5\textwidth}
			\textbf{Extract the secret:}
			\begin{align*}
				g^r                & = Y \cdot A^c                \\
				g^{r'}             & = Y \cdot A^{c'}             \\
				\frac{g^r}{g^{r'}} & = \frac{A^c}{A^{c'}}         \\
				g^{r-r'}           & = A^{c-c'} = g^{a(c-c')}     \\
				r-r'               & = a(c-c') \pmod{n}           \\
				a                  & = \frac{r-r'}{c-c'} \pmod{n}
			\end{align*}
		\end{column}
	\end{columns}
\end{frame}

\begin{frame}{Your turn: extract a mystery key}
	\begin{columns}[T]
		\begin{column}{0.48\textwidth}
			Bob's public key is $B = 3$ (same toy group). You recorded two of his accepting transcripts:
			\begin{center}
				$(Y, c, r) = (4, 1, 10)$ \qquad $(Y, c', r') = (4, 4, 1)$
			\end{center}
			\begin{enumerate}
				\item Check that both transcripts pass
				\item Recover Bob's private key $b$
			\end{enumerate}
			\textbf{Hint:} $1/8 \equiv 7 \pmod{11}$, since $8 \cdot 7 = 56 \equiv 1$.
		\end{column}
		\begin{column}{0.48\textwidth}
			\pause
			\textbf{1. Both pass} ($4 = 2^2$ and $3 = 2^8$):
			\begin{itemize}
				\item $Y \cdot B^{1} = 2^{2 + 8} = 2^{10} = 2^{r}$ \mycheckmark
				\item $Y \cdot B^{4} = 2^{2 + 32} = 2^{1} = 2^{r'}$ \mycheckmark
			\end{itemize}
			\pause
			\textbf{2. Extract:}
			\[
				b = \frac{10 - 1}{1 - 4} = \frac{9}{-3} = \frac{9}{8} = 9 \cdot 7 = 63 \equiv 8
			\]
			And indeed $2^8 = 3 = B$. Bob reused $Y$, and that was enough.
		\end{column}
	\end{columns}
\end{frame}

\begin{frame}{Why extraction implies soundness}
	\begin{columns}[c]
		\begin{column}{0.5\textwidth}
			\textbf{The contrapositive argument:}
			\begin{itemize}
				\item Suppose a prover doesn't know the witness
				\item Then for each $Y$ it sends, it can answer at most \textit{one} challenge (two answers would let us extract!)
				\item The challenge is random, chosen after $Y$
				\item So it passes with probability at most $1 / |\mathcal{C}|$
			\end{itemize}
		\end{column}
		\begin{column}{0.5\textwidth}
			\begin{alertblock}{Soundness guarantee}
				Cheating probability $\leq \frac{1}{\text{number of challenges}}$
			\end{alertblock}
			\begin{exampleblock}{For Schnorr}
				$|\mathcal{C}| = n$, a 256-bit prime\\
				Cheating probability $\approx 2^{-256}$\\
				(The cave: $|\mathcal{C}| = 2$, hence the many rounds.)
			\end{exampleblock}
		\end{column}
	\end{columns}
\end{frame}

\begin{frame}{The extractor, in the wild: never reuse $y$}
	\begin{columns}[c]
		\begin{column}{0.5\textwidth}
			\textbf{The extractor is a thought experiment\ldots}
			\begin{itemize}
				\item \ldots until a real prover reuses its mask $y$
				\item Same $y$, same $Y$: two transcripts with one commitment, exactly the extractor's input
				\item Anyone who sees both computes $a = (r - r')/(c - c')$
			\end{itemize}
			\begin{exampleblock}{The fix}
				Derive $y$ from the key and the message (Ed25519; RFC 6979 for ECDSA): no randomness to fail.
			\end{exampleblock}
		\end{column}
		\begin{column}{0.5\textwidth}
			\begin{alertblock}{It keeps happening}
				\begin{itemize}
					\item \textbf{Sony PlayStation 3 (2010):} its ECDSA signatures used the same nonce every time; fail0verflow recovered Sony's signing key
					\item \textbf{Android (2013):} a random number generator bug repeated ECDSA nonces; Bitcoin wallets were emptied
				\end{itemize}
				ECDSA isn't Schnorr, but the algebra of the attack is the same.
			\end{alertblock}
		\end{column}
	\end{columns}
\end{frame}

\begin{frame}{Zero-knowledge: the formal definition}
	\begin{block}{Definition: Zero-Knowledge}
		A Sigma protocol is \textbf{zero-knowledge} if there exists an efficient \textbf{simulator} that, without the witness, outputs transcripts indistinguishable from real ones.
	\end{block}
	\begin{columns}[T]
		\begin{column}{0.5\textwidth}
			\textbf{Why this captures ``learns nothing'':}
			\begin{itemize}
				\item Whatever the verifier could compute from a real transcript\ldots
				\item \ldots it could compute from a simulated one
				\item And simulated ones contain no information about the witness
			\end{itemize}
		\end{column}
		\begin{column}{0.5\textwidth}
			\textbf{We've seen it three times:}
			\begin{itemize}
				\item The cave: Mallory's edited film
				\item Robin Hood: the Sheriff's paint bucket
				\item The proof: $\lib{}{schnorr-real} \equiv \lib{}{schnorr-fake}$
			\end{itemize}
		\end{column}
	\end{columns}
\end{frame}

\begin{frame}{Honest-verifier zero-knowledge}
	\begin{columns}[c]
		\begin{column}{0.5\textwidth}
			\textbf{A subtlety in our simulator:}
			\begin{itemize}
				\item It picks $c$ itself, uniformly and independently of $Y$
				\item So it only matches verifiers that do the same: \textbf{honest-verifier} zero-knowledge (HVZK)
			\end{itemize}
			\textbf{A cheating verifier} could compute $c = \func{hash}{Y}$ instead. ``Pick $c$ first'' no longer works, and the transcript becomes evidence that the prover took part.
		\end{column}
		\begin{column}{0.5\textwidth}
			\begin{exampleblock}{Good news}
				For most applications, HVZK is enough, and there are standard fixes when it isn't (e.g., the verifier commits to $c$ before seeing $Y$)
			\end{exampleblock}
			\begin{alertblock}{Remember $c = \func{hash}{Y}$}
				This cheating verifier's trick is, almost exactly, the next section's big idea: Fiat-Shamir.
			\end{alertblock}
		\end{column}
	\end{columns}
\end{frame}

\section{Proving AND, OR}

\begin{frame}{Combining proofs: AND and OR}
	\begin{columns}[T]
		\begin{column}{0.5\textwidth}
			\textbf{OR:} one condition holds, and you don't say which
			\begin{itemize}
				\item ``I know the private key for $A_1$ OR $A_2$'': \textbf{anonymous authentication}
				\item ``This encrypted vote is YES or NO'': \textbf{e-voting}. Only valid votes count, and nobody learns which
			\end{itemize}
		\end{column}
		\begin{column}{0.5\textwidth}
			\textbf{AND:} every condition holds
			\begin{itemize}
				\item ``I know $a$ AND $b$ such that\ldots''
				\item ``These discrete logs are equal'': e.g., ``I decrypted this ballot with the election's real key''
			\end{itemize}
		\end{column}
	\end{columns}
	\vspace{0.2cm}
	\begin{alertblock}{The challenge}
		Combine Sigma protocols so that the result is still zero-knowledge, and still a Sigma protocol (so we can keep combining).
	\end{alertblock}
\end{frame}

\subsection{OR Proofs}

\begin{frame}{The Robin Hood analogy}
	\begin{columns}[c]
		\begin{column}{0.5\textwidth}
			\textbf{Robin Hood has two bows:}
			\begin{itemize}
				\item One shoots true, one is erratic
				\item Robin must prove he can shoot, but won't say \textit{which} bow is good
			\end{itemize}
			\textbf{The King's test:}
			\begin{itemize}
				\item Robin fires one arrow; the King paints a target; Robin fires the other
				\item The arrows must sit at equal distances from the target, on opposite sides
			\end{itemize}
		\end{column}
		\begin{column}{0.5\textwidth}
			\textbf{Robin's strategy:}
			\begin{enumerate}
				\item Fire the erratic bow first: it lands where it lands
				\item The King paints the target
				\item Work out the spot that balances the first arrow
				\item Hit it with the good bow
			\end{enumerate}
			\begin{alertblock}{The mapping}
				Erratic arrow $=$ simulated transcript; target $=$ challenge $C$; balance $=$ $C_0 + C_1 = C$.
			\end{alertblock}
		\end{column}
	\end{columns}
\end{frame}

\begin{frame}{OR proofs: the construction}
	\textbf{Given:} Sigma protocols $\Sigma_0$ for condition $P_0$ and $\Sigma_1$ for condition $P_1$\\
	\textbf{Goal:} Prove ``I know $W$ such that $P_0(X,W)$ OR $P_1(X,W)$''
	\begin{block}{The protocol}
		\begin{enumerate}
			\item \textbf{Commit:} Prover sends $(Y_0, Y_1)$, commitments for both protocols
			\item \textbf{Challenge:} Verifier sends single challenge $C$
			\item \textbf{Constraint:} Prover must provide $(C_0, R_0)$ and $(C_1, R_1)$ such that:
			      \begin{itemize}
				      \item $C_0 + C_1 = C$ (challenges sum to verifier's challenge)
				      \item Both transcripts $(Y_0, C_0, R_0)$ and $(Y_1, C_1, R_1)$ are valid
			      \end{itemize}
		\end{enumerate}
	\end{block}
	\textbf{Why sound?} Once $C$ arrives, a prover who knows neither witness is stuck with \textit{two} challenges it didn't choose (it can fix at most one of them in advance).
\end{frame}

\begin{frame}{OR proofs: how it works}
	\begin{columns}[c]
		\begin{column}{0.46\textwidth}
			\textbf{If the prover knows a witness for $P_0$:}
			\begin{enumerate}
				\item Simulate $\Sigma_1$ (Mallory's bet!): pick $C_1, R_1$ at random, compute $Y_1$ backwards
				\item Commit honestly for $\Sigma_0$: real $Y_0$
				\item Receive the challenge $C$
				\item Set $C_0 = C - C_1$
				\item Answer $C_0$ with the witness: real $R_0$
			\end{enumerate}
		\end{column}
		\begin{column}{0.52\textwidth}
			\centering
			\begin{tikzpicture}[crypto diagram]
				\node[cd heading] at (2.9,0.35) {Splitting the verifier's challenge};
				\draw[cd card] (0,-0.1) rectangle (5.8,-0.75);
				\node[font=\small\bfseries] at (2.9,-0.43) {$C$ (verifier's challenge)};
				\draw[cd gold] (0,-1.15) rectangle (2.2,-1.8);
				\node[font=\small\bfseries, text=diagramgold] at (1.1,-1.48) {$C_1$};
				\draw[cd green] (2.3,-1.15) rectangle (5.8,-1.8);
				\node[font=\small\bfseries, text=cipherprimary] at (4.05,-1.48) {$C_0 = C - C_1$};
				\node[cd small, text=diagramgold, anchor=north] at (1.1,-1.85) {chosen \textit{first}:\\simulated};
				\node[cd small, text=cipherprimary, anchor=north] at (4.05,-1.85) {whatever is left:\\answered with the witness};
				\node[cd green, text width=5.5cm, font=\fontsize{9}{11}\selectfont, anchor=north] at (2.9,-2.75) {The verifier sees two uniform-looking pieces that add up to $C$. Which one was chosen first? No way to tell, so it can't tell which witness you hold.};
			\end{tikzpicture}
		\end{column}
	\end{columns}
\end{frame}

\begin{frame}{OR proof protocol diagram}
	\vspace{-1.5cm}
	\begin{center}\resizebox{0.74\textwidth}{!}{
			\begin{msc}{}
				\setmscvalues{small}
				\drawframe{none}
				\setlength{\instdist}{8cm}
				\setlength{\instwidth}{5cm}
				\declinst{prover}{}{Prover (knows $W$ for $P_b$)}
				\declinst{verifier}{}{Verifier}
				\action*{$\begin{array}{l}
							\text{Generate real commitment for } \Sigma_b            \\
							\text{Simulate fake transcript for } \Sigma_{1-b}        \\
							(Y_{1-b}, C_{1-b}, R_{1-b}) \leftarrow \text{Simulate}() \\
						\end{array}$}{prover}
				\nextlevel[5]
				\mess{$(Y_0, Y_1)$}{prover}{verifier}
				\nextlevel[1]
				\action*{$C \twoheadleftarrow \mathcal{C}$}{verifier}
				\nextlevel[2]
				\mess{$C$}{verifier}{prover}
				\nextlevel[1]
				\action*{$\begin{array}{l}
							C_b := C - C_{1-b} \\
							R_b \leftarrow \text{Respond using witness}
						\end{array}$}{prover}
				\nextlevel[4]
				\mess{$(C_0, R_0, C_1, R_1)$}{prover}{verifier}
				\nextlevel[1]
				\action*{$\begin{array}{l}
							\text{Check: } C_0 + C_1 \overset{?}{=} C            \\
							\text{Check: } \Sigma_0.\text{Verify}(Y_0, C_0, R_0) \\
							\text{Check: } \Sigma_1.\text{Verify}(Y_1, C_1, R_1)
						\end{array}$}{verifier}
				\nextlevel[4]
			\end{msc}
		}\end{center}
\end{frame}

\begin{frame}{Example: an OR proof with toy numbers}
	\centering
	\begin{tikzpicture}[crypto diagram]
		\path[use as bounding box] (0,0) rectangle (13.5,-5.7);
		\node[anchor=west, font=\fontsize{9.5}{11}\selectfont] at (0,-0.2) {Public keys $A_0 = 13$ and $A_1 = 6$. Alice knows $a_0 = 7$, but not $a_1$. She proves she knows one of them.};
		\node[cd heading] at (5.55,-0.75) {Branch 0: real (knows $a_0 = 7$)};
		\node[cd heading] at (10.6,-0.75) {Branch 1: simulated};
		\foreach \y/\step/\real/\fake in {
				-1.35/{1. Commit}/{random $y_0 = 5$: $Y_0 = 2^5 = 9$}/{pick $C_1 = 4$, $R_1 = 5$: $Y_1 = 2^{5} \cdot (6^{4})^{-1} = 9 \cdot 8^{-1} = 4$},
				-2.2/{2. $C = 6$ arrives}/{$C_0 = C - C_1 = 6 - 4 = 2$}/{already fixed: $C_1 = 4$},
				-3.05/{3. Respond}/{$R_0 = y_0 + C_0 a_0 = 5 + 14 = 19 \equiv 8$}/{already fixed: $R_1 = 5$},
				-3.9/{4. Verifier checks}/{$2^{8} = 3$ and $Y_0 A_0^{C_0} = 9 \cdot 13^{2} = 3$ \mycheckmark}/{$2^{5} = 9$ and $Y_1 A_1^{C_1} = 4 \cdot 6^{4} = 9$ \mycheckmark}
			} {
				\node[anchor=west, font=\fontsize{9}{11}\selectfont\bfseries, text=cipherprimarydark] at (0,\y) {\step};
				\node[cd green, minimum width=4.95cm, minimum height=0.66cm, inner sep=2pt, font=\fontsize{8.5}{10}\selectfont] at (5.55,\y) {\real};
				\node[cd gold, minimum width=4.95cm, minimum height=0.66cm, text width=4.8cm, inner sep=2pt, font=\fontsize{8.5}{10}\selectfont] at (10.6,\y) {\fake};
			}
		\node[cd card, text width=12.9cm, font=\fontsize{9}{11}\selectfont] at (6.75,-4.95) {Also $C_0 + C_1 = 2 + 4 = 6 = C$ \mycheckmark. The verifier gets $(C_0, R_0, C_1, R_1) = (2, 8, 4, 5)$: two random-looking challenge--response pairs that both pass. Nothing marks branch 1 as the fake.};
	\end{tikzpicture}
\end{frame}

\subsection{AND Proofs}

\begin{frame}{AND proofs: one challenge for everyone}
	\begin{columns}[c]
		\begin{column}{0.42\textwidth}
			\textbf{The problem:}
			\begin{itemize}
				\item Prove that several conditions hold at once
				\item Running the protocols separately gives several challenges: no longer a single Sigma protocol, so we can't nest it inside an OR
			\end{itemize}
			\textbf{The fix:} commit to everything, then answer \textit{one} shared challenge. Soundness and zero-knowledge both carry over.
		\end{column}
		\begin{column}{0.56\textwidth}
			\begin{block}{The protocol, for $\Sigma_1, \ldots, \Sigma_k$}
				\begin{enumerate}
					\item \textbf{Commit:} $(Y_i, \sigma_i) \leftarrow \Sigma_i.\text{Commit}(X_i, W_i)$ for all $i$; send $(Y_1, \ldots, Y_k)$
					\item \textbf{Challenge:} receive a single $C$
					\item \textbf{Respond:} $R_i \leftarrow \Sigma_i.\text{Respond}(\sigma_i, C)$ for all $i$; send $(R_1, \ldots, R_k)$
					\item \textbf{Verify:} every $(Y_i, C, R_i)$ is valid
				\end{enumerate}
			\end{block}
		\end{column}
	\end{columns}
\end{frame}

\begin{frame}{Example: proving equality of discrete logs}
	\textbf{Setup:} $(g_1, g_2, g_3, A_1, A_2, A_3)$ with $A_i = g_i^a$ for the same $a$; prove knowledge of $a$ (Chaum-Pedersen, 1992):
	\begin{columns}[c]
		\begin{column}{0.5\textwidth}
			\begin{enumerate}
				\item \textbf{Commit:}
				      \begin{itemize}
					      \item Pick one $y \twoheadleftarrow \mathbb{Z}_n$
					      \item Send $Y_i = g_i^y$ for $i = 1, 2, 3$
				      \end{itemize}
				\item \textbf{Challenge:} Receive $C$
				\item \textbf{Respond:} Send $R = y + C \cdot a$
			\end{enumerate}
		\end{column}
		\begin{column}{0.5\textwidth}
			\begin{enumerate}
				\setcounter{enumi}{3}
				\item \textbf{Verify:}
				      \begin{itemize}
					      \item Check $g_1^R = Y_1 \cdot A_1^C$
					      \item Check $g_2^R = Y_2 \cdot A_2^C$
					      \item Check $g_3^R = Y_3 \cdot A_3^C$
				      \end{itemize}
			\end{enumerate}
		\end{column}
	\end{columns}
	\begin{alertblock}{One $y$, one $R$: that is what proves equality}
		Separate runs allow different $a$'s; one shared $y$ and $R$ forces the \textit{same} $a$ everywhere. You'll meet it in e-voting (proving a ballot was decrypted with the real key) and in Privacy Pass (a server proving it used the same key for every user).
	\end{alertblock}
\end{frame}

\begin{frame}{Combining AND and OR}
	\begin{columns}[c]
		\begin{column}{0.45\textwidth}
			\textbf{Example:} ``I know the private key for ($A_1$ OR $A_2$) AND ($B_1$ OR $B_2$)''
			\begin{itemize}
				\item Four possible combinations; the verifier learns nothing about which
			\end{itemize}
			\textbf{Construction:}
			\begin{enumerate}
				\item An OR proof for each clause
				\item Combine them with AND
				\item Result: one Sigma protocol for the whole formula, so you can keep nesting
			\end{enumerate}
		\end{column}
		\begin{column}{0.53\textwidth}
			\centering
			\begin{tikzpicture}[crypto diagram, level distance=1.35cm]
				\node[cd card, font=\small\bfseries, minimum width=1.6cm] (and) at (3.3,0) {AND};
				\node[cd card, font=\small\bfseries, minimum width=1.4cm] (or1) at (1.55,-1.45) {OR};
				\node[cd card, font=\small\bfseries, minimum width=1.4cm] (or2) at (5.05,-1.45) {OR};
				\node[cd green, minimum width=1.25cm] (a1) at (0.7,-2.95) {$A_1$};
				\node[cd gold, minimum width=1.25cm] (a2) at (2.4,-2.95) {$A_2$};
				\node[cd gold, minimum width=1.25cm] (b1) at (4.2,-2.95) {$B_1$};
				\node[cd green, minimum width=1.25cm] (b2) at (5.9,-2.95) {$B_2$};
				\draw[cd arrow, -] (and) -- node[cd tag, pos=0.45] {$C$} (or1);
				\draw[cd arrow, -] (and) -- node[cd tag, pos=0.45] {$C$} (or2);
				\draw[cd arrow, -] (or1) -- (a1);
				\draw[cd arrow, -, draw=diagramgold] (or1) -- (a2);
				\draw[cd arrow, -, draw=diagramgold] (or2) -- (b1);
				\draw[cd arrow, -] (or2) -- (b2);
				\node[cd small, text=cipherdark!70] at (1.55,-3.65) {$C_{A_1} + C_{A_2} = C$};
				\node[cd small, text=cipherdark!70] at (5.05,-3.65) {$C_{B_1} + C_{B_2} = C$};
				\node[cd small, text width=6cm] at (3.3,-4.45) {Knowing $a_1$ and $b_2$: two \textcolor{cipherprimary}{\textbf{real}} leaves, two \textcolor{diagramgold}{\textbf{simulated}}. All four transcripts look alike.};
			\end{tikzpicture}
		\end{column}
	\end{columns}
\end{frame}

\section{Non-Interactive Zero-Knowledge Proofs}

\subsection{Fiat-Shamir}

\begin{frame}{Non-interactive proofs: removing the verifier}
	\begin{center}
		\textbf{The limitation of interactive proofs}
	\end{center}
	\begin{columns}[c]
		\begin{column}{0.5\textwidth}
			\textbf{Interactive Sigma protocols:}
			\begin{itemize}
				\item Require back-and-forth communication
				\item Prover and verifier must be online together
				\item Can't post a proof for later verification
				\item Can't broadcast a proof to many verifiers
			\end{itemize}
		\end{column}
		\begin{column}{0.5\textwidth}
			\textbf{What we want:}
			\begin{itemize}
				\item Prover generates proof alone
				\item Anyone can verify later
				\item No interaction needed
				\item But still zero-knowledge!
			\end{itemize}
			\begin{alertblock}{The challenge}
				Without a verifier, who guarantees that $C$ is unpredictable and comes \textit{after} $Y$?
			\end{alertblock}
		\end{column}
	\end{columns}
\end{frame}

\begin{frame}{The Fiat-Shamir transformation (1986)}
	\centering
	\begin{tikzpicture}[crypto diagram]
		\path[use as bounding box] (0,0) rectangle (13.5,-5.7);
		% Interactive
		\node[cd heading] at (2.9,-0.2) {Interactive};
		\node[cd green, minimum width=1.7cm] (p1) at (0.95,-0.85) {\textbf{Prover}};
		\node[cd card, minimum width=1.7cm] (v1) at (4.85,-0.85) {\textbf{Verifier}};
		\draw[ciphergraydark, line width=1.2pt] (0.95,-1.2) -- (0.95,-3.55);
		\draw[ciphergraydark, line width=1.2pt] (4.85,-1.2) -- (4.85,-3.55);
		\draw[cd arrow] (0.95,-1.65) -- node[cd tag, above=1pt] {$Y$} (4.85,-1.65);
		\draw[cd arrow, draw=diagramgold] (4.85,-2.45) -- node[cd tag, above=1pt] {random $c$} (0.95,-2.45);
		\draw[cd arrow] (0.95,-3.25) -- node[cd tag, above=1pt] {$r$} (4.85,-3.25);
		\draw[diagramgold!60, line width=1pt, dashed] (6.45,-0.4) -- (6.45,-3.8);
		% Non-interactive
		\node[cd heading] at (10.05,-0.2) {Fiat-Shamir: the hash plays the verifier};
		\node[cd green, minimum width=1.7cm] at (7.9,-0.85) {\textbf{Prover}};
		\node[cd card, minimum width=1.9cm, font=\small] (y) at (7.9,-1.65) {$Y = g^y$};
		\node[cd gold, minimum width=2.4cm, font=\small] (h) at (10.8,-1.65) {$c = \func{hash}{X, Y}$};
		\node[cd card, minimum width=1.9cm, font=\small] (r) at (7.9,-2.55) {$r = y + ca$};
		\draw[cd arrow] (y) -- (h);
		\draw[cd arrow, draw=diagramgold] (h.south) |- (r.east);
		\draw[cd arrow] (y) -- (r);
		\node[cd green, minimum width=2.6cm, font=\small] (proof) at (8.55,-3.55) {proof $\pi = (Y, r)$};
		\draw[cd arrow] (r) -- (7.9,-3.2);
		\node[cd small, text=cipherdark!70, anchor=west] at (10.0,-3.55) {anyone, any time: recompute\\$c = \func{hash}{X, Y}$, then check};
		% Random oracle
		\node[cd gold, text width=12.9cm, align=flush left, font=\fontsize{9}{11}\selectfont] at (6.75,-5.05) {\textbf{The model:} treat the hash as a \textbf{random oracle}, a fresh uniform output for every new input and the same output for a repeated one. SHA-256 stands in for it in practice. Since $c$ depends on $Y$, the prover must fix $Y$ first, just as with a live verifier.};
	\end{tikzpicture}
\end{frame}

\begin{frame}{Fiat-Shamir: the construction}
	Given a Sigma protocol $\Sigma$, create a non-interactive proof system:\\[1em]
	\begin{columns}[c]
		\begin{column}{0.5\textwidth}
			\textbf{Prove}$(X, W)$:
			\begin{enumerate}
				\item $(Y, \sigma) := \Sigma.\text{Commit}(X, W)$
				\item \hl{$C := \func{hash}{X, Y}$}
				\item $R := \Sigma.\text{Respond}(\sigma, C)$
				\item Return $(Y, R)$
			\end{enumerate}
		\end{column}
		\begin{column}{0.5\textwidth}
			\textbf{Verify}$(X, Y, R)$:
			\begin{enumerate}
				\item \hl{$C := \func{hash}{X, Y}$}
				\item Return $\Sigma.\text{Check}(X, Y, C, R)$
			\end{enumerate}
		\end{column}
	\end{columns}
	\begin{center}
		\textbf{The proof is just $(Y, R)$: no interaction needed!}
	\end{center}
\end{frame}

\begin{frame}{Example: non-interactive Schnorr}
	\begin{columns}[c]
		\begin{column}{0.5\textwidth}
			\textbf{Interactive Schnorr:}
			\begin{enumerate}
				\item Prover sends $Y = g^y$
				\item Verifier sends random $c$
				\item Prover sends $r = y + ca$
				\item Verifier checks $g^r = Y \cdot A^c$
			\end{enumerate}
		\end{column}
		\begin{column}{0.5\textwidth}
			\textbf{Fiat-Shamir Schnorr:}
			\begin{enumerate}
				\item Prover computes $Y = g^y$
				\item Prover computes $c = \func{hash}{A, Y}$
				\item Prover computes $r = y + ca$
				\item Proof is $(Y, r)$
				\item Verifier recomputes $c = \func{hash}{A, Y}$
				\item Verifier checks $g^r = Y \cdot A^c$
			\end{enumerate}
		\end{column}
	\end{columns}
	\begin{alertblock}{This is a digital signature!}
		Put a message $m$ in the hash, $c = \func{hash}{A, Y, m}$, and you get the Schnorr signature scheme: a proof of knowledge of $a$ that is bound to $m$.
	\end{alertblock}
\end{frame}

\begin{frame}{Why Fiat-Shamir is sound}
	\begin{columns}[c]
		\begin{column}{0.5\textwidth}
			\textbf{The hash enforces the order:}
			\begin{itemize}
				\item The prover must fix $Y$ first
				\item Only then does $C = \func{hash}{X, Y}$ exist
				\item In the random oracle model, $C$ is unpredictable until $Y$ is fixed, just like a verifier's challenge
				\item So the prover can't ``work backwards'' from $C$ to $Y$ the way our simulator did
			\end{itemize}
		\end{column}
		\begin{column}{0.5\textwidth}
			\begin{alertblock}{Still zero-knowledge?}
				Yes, in the random oracle model: the simulator picks $C, R$, computes $Y$ backwards, then \textit{programs} the oracle so $\func{hash}{X, Y} := C$.
			\end{alertblock}
			\begin{exampleblock}{Why hash $X$ too?}
				Omit the statement (``weak'' Fiat-Shamir) and a prover can fix $Y$ first, then pick $X$. This broke Helios's voting proofs (Bernhard-Pereira-Warinschi, 2012).
			\end{exampleblock}
		\end{column}
	\end{columns}
\end{frame}

\begin{frame}{Fiat-Shamir needs a huge challenge space}
	\begin{columns}[c]
		\begin{column}{0.47\textwidth}
			\textbf{Make the cave non-interactive:} Peggy's challenge bit is now $\func{hash}{\text{her commitment}}$.
			\begin{itemize}
				\item A cheater just retries: commit, hash, and if the bit names the wrong side, start over
				\item Two tries on average, and no one sees the failures
				\item Live, Victor would never let her retry. Offline, nobody is watching
			\end{itemize}
		\end{column}
		\begin{column}{0.5\textwidth}
			\begin{center}
				\small
				\begin{tabular}{lll}
					\textbf{Challenge} & \textbf{Live: fooled} & \textbf{FS: cheat costs} \\[0.3em]
					1 bit              & $1/2$                 & 2 hashes                 \\
					20 bits            & $1$ in a million      & $10^6$ hashes: $<$1 s  \\
					128 bits           & $2^{-128}$            & $2^{128}$ hashes: never
				\end{tabular}
			\end{center}
			\begin{alertblock}{The rule}
				Fiat-Shamir soundness is counted in hash evaluations. Schnorr's 256-bit $c$ is fine; protocols with small challenges need enough parallel repetitions to reach 128 bits.
			\end{alertblock}
		\end{column}
	\end{columns}
\end{frame}

\begin{frame}{The forking lemma: rewind and fork}
	\centering
	\begin{tikzpicture}[crypto diagram]
		\path[use as bounding box] (0,0) rectangle (13.5,-5.7);
		\node[anchor=west, font=\fontsize{9.5}{11}\selectfont, text width=13.3cm, align=flush left] at (0,-0.35) {\textbf{Forking lemma} (Pointcheval-Stern, 1996): if $\mathcal{A}$ outputs valid Fiat-Shamir proofs with non-negligible probability, rewinding it as below yields a witness with non-negligible probability.};
		\node[cd card, minimum width=1.6cm] (start) at (0.85,-2.2) {run $\mathcal{A}$};
		\node[cd green, minimum width=2.6cm, font=\fontsize{9}{11}\selectfont] (query) at (3.6,-2.2) {$\mathcal{A}$ asks for\\$\func{hash}{X, Y}$};
		\draw[cd arrow] (start) -- (query);
		\node[cd card, minimum width=3.3cm, font=\fontsize{9}{11}\selectfont] (up) at (7.6,-1.45) {answer $C$ $\Rightarrow$ proof $(Y, R)$};
		\node[cd gold, minimum width=3.3cm, font=\fontsize{9}{11}\selectfont] (down) at (7.6,-2.95) {rewind, answer $C'$\\$\Rightarrow$ proof $(Y, R')$};
		\draw[cd arrow] (query.east) -- ++(0.35,0) |- (up.west);
		\draw[cd arrow, draw=diagramgold] (query.east) -- ++(0.35,0) |- (down.west);
		\node[cd tag, text=diagramgold] at (5.2,-2.2) {fork};
		\node[cd green, minimum width=2.9cm, font=\fontsize{9}{11}\selectfont] (ext) at (11.75,-2.2) {\textbf{Extract}\\$(Y, C, R, C', R')$\\$\Rightarrow W$};
		\draw[cd arrow] (up.east) -| (ext.north);
		\draw[cd arrow] (down.east) -| (ext.south);
		\node[cd card, text width=6.1cm, align=flush left, font=\fontsize{8.5}{10.5}\selectfont, anchor=north west] at (0,-3.8) {\textbf{Same $Y$ on both branches:} everything before the fork replays identically, same random coins included. Special soundness does the rest.};
		\node[cd gold, text width=6.3cm, align=flush left, font=\fontsize{8.5}{10.5}\selectfont, anchor=north east] at (13.5,-3.8) {\textbf{But hashes are deterministic!} The extractor \textit{simulates} the random oracle, so it may answer the same query differently after rewinding. A thought experiment: you can't do this to SHA-256.};
	\end{tikzpicture}
\end{frame}

\begin{frame}{Deniability lost!}
	\begin{columns}[c]
		\begin{column}{0.5\textwidth}
			\textbf{Interactive Sigma protocol:}
			\begin{itemize}
				\item Transcripts are deniable
				\item Anyone can fake one: choose $C, R$ first, compute $Y$
				\item Indistinguishable from real
				\item \textbf{Can't prove anything to third parties}
			\end{itemize}
		\end{column}
		\begin{column}{0.5\textwidth}
			\textbf{Fiat-Shamir proof:}
			\begin{itemize}
				\item $C = \func{hash}{X, Y}$, checkable by everyone
				\item ``Choose $C$ first'' fails: $C$ must hash a $Y$ you haven't computed yet
				\item Only the witness holder can create one
				\item \textbf{Transferable evidence!}
			\end{itemize}
		\end{column}
	\end{columns}
	\begin{alertblock}{Trade-off}
		We gain non-interactivity but lose deniability: Victor's film now has the hash baked in, so Mallory can't stage it. Exactly what a signature needs.
	\end{alertblock}
\end{frame}

\begin{frame}[fragile]{Your turn: spot the bugs}
	\begin{columns}[T]
		\begin{column}{0.46\textwidth}
			Log in by proving you own public key \texttt{A}:
			{\small
\begin{verbatim}
def prove(a):
    y = random_scalar()
    Y = g ** y
    c = H(Y)
    r = (y + c * a) % n
    return (Y, r)

def verify(A, Y, r):
    c = H(Y)
    return g ** r == Y * A ** c
\end{verbatim}
			}
			SHA-256, 256-bit $n$. What's wrong?
		\end{column}
		\begin{column}{0.52\textwidth}
			\pause
			\textbf{Bug 1: the statement isn't hashed.} Mallory takes someone else's key $B$ as her ``commitment'': $Y = B$, $c = \func{hash}{B}$, picks any $r$, then solves for $A = (g^r / B)^{1/c}$. A valid proof for a key whose secret she doesn't know. Fix: $c = \func{hash}{A, Y}$.
			\pause
			\vspace{0.2cm}

			\textbf{Bug 2: no context is hashed.} The proof is transferable, so an eavesdropper who sees one login can replay $(Y, r)$ forever. Fix: hash a fresh server nonce and the session: $c = \func{hash}{A, Y, \text{nonce}, \text{session}}$.
		\end{column}
	\end{columns}
\end{frame}

\begin{frame}{Fiat-Shamir is everywhere}
	\begin{columns}[c]
		\begin{column}{0.5\textwidth}
			\textbf{Signatures you already use:}
			\begin{itemize}
				\item Schnorr signatures (in Bitcoin since Taproot, 2021)
				\item Ed25519 (SSH, TLS, Signal): Schnorr over Curve25519
				\item ML-DSA (Dilithium), from lecture 2.6: ``Fiat-Shamir with aborts''
				\item ECDSA is the odd one out: not a Fiat-Shamir construction
			\end{itemize}
		\end{column}
		\begin{column}{0.5\textwidth}
			\textbf{Proof systems:}
			\begin{itemize}
				\item Bulletproofs: the range proofs hiding Monero amounts
				\item Every SNARK and STARK: an interactive protocol, made non-interactive by hashing
			\end{itemize}
			\begin{alertblock}{Keep this in mind}
				All of these lean on the random oracle model. Next: what happens when the hash is a real function and the protocol is a big one.
			\end{alertblock}
		\end{column}
	\end{columns}
\end{frame}

\subsection{From Statements to Circuits}

\begin{frame}{Why circuits? One protocol for every statement}
	\begin{center}
		\textbf{Claim: ``I know $x$ such that SHA-256($x$) = $y$''}
	\end{center}
	\begin{columns}[c]
		\begin{column}{0.5\textwidth}
			\textbf{As a hand-made Sigma protocol:}
			\begin{itemize}
				\item SHA-256 has 64 rounds of bitwise operations
				\item No algebraic structure to exploit, unlike $g^a$
				\item Only the generic NP route remains: hopelessly impractical
			\end{itemize}
		\end{column}
		\begin{column}{0.5\textwidth}
			\textbf{As a circuit:}
			\begin{itemize}
				\item Write SHA-256 as a circuit: private input $x$, public output $y$
				\item Run one general-purpose circuit ZK protocol
				\item Done, at a price: tens of thousands of constraints for a single block
			\end{itemize}
		\end{column}
	\end{columns}
	\begin{alertblock}{The shift}
		Old way: a clever protocol per statement. New way: write the computation as a circuit, and use one protocol for everything: ``I know $w$ such that $C(x, w) = y$''.
	\end{alertblock}
\end{frame}

\begin{frame}{What is a circuit? A worked example}
	\centering
	\begin{tikzpicture}[crypto diagram]
		\path[use as bounding box] (0,0) rectangle (13.5,-5.7);
		\node[anchor=west, font=\fontsize{9.5}{11}\selectfont] at (0,-0.25) {\textbf{Claim:} ``I know $x$ such that $x^3 + x + 5 = 35$.'' The witness is $x = 3$.};
		\node[cd gold, circle, minimum size=0.85cm, font=\small\bfseries] (x) at (0.6,-2.6) {$x$};
		\node[cd small, text=diagramgold] at (0.6,-3.3) {private};
		\node[cd green, minimum width=1cm, minimum height=0.8cm, font=\bfseries] (m1) at (2.3,-1.5) {$\times$};
		\node[cd green, minimum width=1cm, minimum height=0.8cm, font=\bfseries] (m2) at (4.3,-1.5) {$\times$};
		\node[cd green, minimum width=1cm, minimum height=0.8cm, font=\bfseries] (a1) at (6.2,-2.6) {$+$};
		\node[cd green, minimum width=1cm, minimum height=0.8cm, font=\bfseries] (a2) at (8.0,-2.6) {$+$};
		\node[cd card, circle, minimum size=0.75cm, font=\small\bfseries] (five) at (8.0,-3.95) {$5$};
		\node[cd small, text=cipherdark!65, anchor=north] at (8.0,-4.35) {public constant};
		\node[cd card, minimum width=1.1cm, minimum height=0.8cm, font=\small\bfseries] (out) at (9.6,-2.6) {$35$};
		\node[cd small, text=cipherdark!65, anchor=north] at (9.6,-3.05) {public\\output};
		\draw[cd arrow] (x.north) |- ([yshift=4pt]m1.west);
		\draw[cd arrow] (x.north east) -- ([yshift=-5pt]m1.west);
		\draw[cd arrow] (m1) -- node[cd tag, text=diagramgold, font=\fontsize{8}{9}\selectfont\bfseries, above=2pt] {$v_1 = 9$} (m2);
		\draw[cd arrow] (x.east) -- (3.3,-2.6) -- (3.3,-1.75) -- ([yshift=-5pt]m2.west);
		\draw[cd arrow] (m2.east) -| node[cd tag, text=diagramgold, font=\fontsize{8}{9}\selectfont\bfseries, pos=0.3, above=2pt] {$v_2 = 27$} (a1.north);
		\draw[cd arrow] (x.south east) -- (1.2,-3.0) -- (5.4,-3.0) -- ([yshift=-5pt]a1.west);
		\draw[cd arrow] (a1) -- node[cd tag, text=diagramgold, font=\fontsize{8}{9}\selectfont\bfseries, above=2pt] {$v_3 = 30$} (a2);
		\draw[cd arrow] (five) -- (a2);
		\draw[cd arrow] (a2) -- (out);
		\node[cd card, text width=5.2cm, align=flush left, font=\fontsize{8.5}{10.5}\selectfont, anchor=north west] at (0,-3.95) {\textbf{Gates:} $+$ and $\times$ modulo a prime. No loops: the circuit has a fixed size, so loops get unrolled.};
		\node[cd green, text width=2.9cm, align=flush left, font=\fontsize{8.5}{10.5}\selectfont, anchor=north east] at (13.5,-0.75) {\textbf{The proof shows each gate holds:}\\$v_1 = x \cdot x$\\$v_2 = v_1 \cdot x$\\$v_3 = v_2 + x$\\$v_3 + 5 = 35$\\\ldots while hiding $x$ and all the \textcolor{diagramgold}{wire values}.};
	\end{tikzpicture}
\end{frame}

\begin{frame}[fragile]{From theory to practice: ZK programming}
	\begin{columns}[c]
		\begin{column}{0.5\textwidth}
			\textbf{Circom's \texttt{IsZero}:} out $= 1$ if in $= 0$, else $0$
			{\small
\begin{verbatim}
template IsZero() {
  signal input in;
  signal output out;
  signal inv;
  inv <-- in != 0 ? 1/in : 0;
  out <== -in*inv + 1;
  in*out === 0;
}
\end{verbatim}
			}
			A compiler turns this into a constraint system; a general proof system proves the constraints hold.
		\end{column}
		\begin{column}{0.5\textwidth}
			\textbf{Read the constraints:}
			\begin{itemize}
				\item \texttt{<--} is a \textit{hint}: the prover just computes \texttt{inv}, nothing is checked
				\item \texttt{<==} and \texttt{===} are \textit{constraints} the proof enforces
				\item If in $\neq 0$: out $= 1 - \text{in} \cdot \text{in}^{-1} = 0$, and in $\cdot$ out $= 0$ holds
				\item If in $= 0$: out $= 1$, and $0 \cdot 1 = 0$ holds
				\item No branches, no loops, only equations: getting here from ordinary code is the compiler's job
			\end{itemize}
		\end{column}
	\end{columns}
\end{frame}

\begin{frame}{Your turn: delete one line}
	\begin{columns}[T]
		\begin{column}{0.48\textwidth}
			Remove the last constraint from \texttt{IsZero}, \texttt{in*out === 0}. The circuit still compiles, and honest provers still get the right answer.
			\vspace{0.2cm}

			\textbf{What can a cheating prover prove now?}
			\pause
			\vspace{0.2cm}

			Take in $= 5$. The hint is unchecked, so set inv $= 0$. Then out $= -5 \cdot 0 + 1 = 1$, and no remaining constraint objects: a valid proof that $5$ is zero.
		\end{column}
		\begin{column}{0.48\textwidth}
			\pause
			\begin{alertblock}{Under-constrained circuits}
				The proof system was fine; the circuit said less than its author meant. One of the most common real ZK bugs.
			\end{alertblock}
			\begin{exampleblock}{It happened: Tornado Cash (2019)}
				A circomlib hash used \texttt{=} where it needed \texttt{<==}: a missing constraint, so anyone could forge withdrawals. The team exploited it first to rescue the funds.
			\end{exampleblock}
		\end{column}
	\end{columns}
\end{frame}

\begin{frame}{The cost of generality}
	\begin{center}
		\textbf{Circuit ZK is powerful but comes with trade-offs}
	\end{center}
	\begin{columns}[c]
		\begin{column}{0.5\textwidth}
			\textbf{Advantages:}
			\begin{itemize}
				\item Universal: works for any computation
				\item Compositional: combine circuits
				\item Tool support: compilers, libraries
				\item No need to design a new protocol
			\end{itemize}
		\end{column}
		\begin{column}{0.5\textwidth}
			\textbf{Disadvantages:}
			\begin{itemize}
				\item Less efficient than custom protocols
				\item Circuit can be huge (millions of gates)
				\item Proving time can be slow
				\item New bug classes: under-constrained circuits
			\end{itemize}
		\end{column}
	\end{columns}
	\begin{exampleblock}{Rule of thumb}
		Use custom Sigma protocols for simple statements (like Schnorr)\\
		Use circuit ZK for complex computations (like program execution)
	\end{exampleblock}
\end{frame}

\subsection{Attacks on Fiat-Shamir}

\begin{frame}{The problem with Fiat-Shamir}
	\begin{center}
		\textbf{Can we actually trust Fiat-Shamir in practice?}
	\end{center}
	\begin{columns}[c]
		\begin{column}{0.5\textwidth}
			\textbf{The promise:}
			\begin{itemize}
				\item FS is secure in random oracle model
				\item Replace oracle with concrete hash
				\item Should be secure in practice
			\end{itemize}
			\textbf{Known theoretical attacks:}
			\begin{itemize}
				\item Canetti-Goldreich-Halevi (2004)
				\item Barak (2001), Goldwasser-Kalai (2003)
				\item But all use contrived protocols!
			\end{itemize}
		\end{column}
		\begin{column}{0.5\textwidth}
			\begin{alertblock}{New result: Khovratovich, Rothblum, Soukhanov (2025)}
				First practical attack on a standard protocol!\footnote{\url{https://appliedcryptography.page/paper/\#prove-false}}
				\begin{itemize}
					\item Targets GKR-based proof systems
					\item Used in real deployed systems
					\item Attack works for \textit{any} hash function
				\end{itemize}
			\end{alertblock}
		\end{column}
	\end{columns}
\end{frame}

\begin{frame}{The GKR protocol under attack}
	\begin{columns}[c]
		\begin{column}{0.5\textwidth}
			\textbf{What is GKR?}
			\begin{itemize}
				\item Protocol by Goldwasser-Kalai-Rothblum
				\item Proves ``$C(x, w) = y$'' for bounded-depth circuits, layer by layer
				\item Combined with a polynomial commitment to $w$
				\item Used in practice (e.g., Expander), and a building block of many SNARKs
			\end{itemize}
		\end{column}
		\begin{column}{0.5\textwidth}
			\begin{exampleblock}{The protocol flow}
				\begin{enumerate}
					\item Commit to the witness $w$
					\item Verifier picks a random point $r$
					\item Compare the claimed output with the real one \textit{at the point $r$}
					\item Run the GKR reduction, check with the polynomial commitment
				\end{enumerate}
			\end{exampleblock}
			With Fiat-Shamir, $r$ is a hash of the circuit, the claim and the commitment.
		\end{column}
	\end{columns}
\end{frame}

\begin{frame}{Understanding the attack: an exam analogy}
	\textbf{The honest exam:}
	\begin{itemize}
		\item You hand in an essay ($C$), a sealed answer sheet ($w$), and a claimed grade ($y^*$)
		\item The professor hashes all three to pick the questions ($r$), so you can't predict them
	\end{itemize}
	\textbf{The cheat:}
	\begin{itemize}
		\item Your essay says: ``hash the ID on my answer sheet with the grade and the sheet; that's the question. Tune the output so the check passes exactly there''
		\item On the answer sheet you write the essay's real ID, $\langle C^* \rangle$: the essay now knows the question before the professor does
	\end{itemize}
	\begin{alertblock}{Why the detour through the answer sheet?}
		The ID is a hash \textit{of the essay}, so the essay can't contain it. Passing it in as the witness breaks the circle.
	\end{alertblock}
\end{frame}

\begin{frame}{Circuit digests: how a circuit can talk about itself}
	\begin{columns}[c]
		\begin{column}{0.5\textwidth}
			\textbf{From circuit to digest:}
			\begin{itemize}
				\item Serialize a circuit $C$ (gates, wires, inputs, outputs) in a fixed order and hash it: the \textbf{circuit digest} $\psi = \langle C \rangle$
				\item $\psi$ goes into the Fiat-Shamir hash: the challenge depends on \textit{which} circuit is proved
			\end{itemize}
			\textbf{The circular dependency:}
			\begin{itemize}
				\item The attacker wants $C^*$ to compute its own challenge $r$
				\item But $r$ depends on $\langle C^* \rangle$, and changing $C^*$ changes $\langle C^* \rangle$
			\end{itemize}
		\end{column}
		\begin{column}{0.5\textwidth}
			\begin{alertblock}{Breaking the cycle}
				\begin{itemize}
					\item A circuit can't contain its own digest
					\item But it can receive it as the witness: $w = \langle C^* \rangle$
					\item Then it can hash its own description, exactly as the verifier will
				\end{itemize}
			\end{alertblock}
			If the circuit is deep enough to evaluate the hash, it can \textbf{predict its own challenge}.
		\end{column}
	\end{columns}
\end{frame}

\begin{frame}{The attack, step by step}
	\centering
	\begin{tikzpicture}[crypto diagram]
		\path[use as bounding box] (0,0) rectangle (13.5,-5.7);
		\draw[cd card, fill=cipherprimary!3!cipherwhite] (0,-0.1) rectangle (5.6,-5.6);
		\node[cd heading] at (2.8,-0.4) {The attacking circuit $C^*$};
		\foreach \y/\name/\content in {
				-1.05/w/{Input: witness $w$},
				-1.9/psi/{Read $w$ as a digest $\psi$},
				-2.75/alpha/{$\alpha = \text{comm}(w)$},
				-3.6/gamma/{$\gamma = \func{hash}{\psi, y^*, \alpha}$},
				-4.45/out/{Output $(\gamma,\ \gamma - 1)$}
			} {
				\node[cd green, minimum width=4.2cm, minimum height=0.6cm, inner sep=2pt, font=\fontsize{9}{11}\selectfont] (\name) at (2.8,\y) {\content};
			}
		\foreach \a/\b in {w/psi,psi/alpha,alpha/gamma,gamma/out}
		\draw[cd arrow] (\a) -- (\b);
		\node[cd small, text=diagramred] at (2.8,-5.1) {Never $(0, 0)$: the claim is \textbf{false}};
		% Prover and verifier
		\node[cd gold, text width=6.8cm, align=flush left, font=\fontsize{9}{11}\selectfont] (claim) at (10,-0.55) {\textbf{Prover:} claims $C^*(w) = y^* = (0, 0)$, and sets the witness to $w = \langle C^* \rangle$};
		\node[cd card, text width=6.8cm, align=flush left, font=\fontsize{9}{11}\selectfont] (ver) at (10,-1.65) {\textbf{Verifier (Fiat-Shamir):} $r = \func{hash}{\langle C^* \rangle, y^*, \text{comm}(w)}$};
		\draw[diagramgold, line width=1.2pt, dashed] (gamma.east) -- (5.95,-3.6) -- (5.95,-1.65) -- (ver.west);
		\node[cd tag, text=diagramgold, font=\fontsize{8.5}{10}\selectfont\bfseries] at (5.95,-2.65) {$r = \gamma$};
		\node[cd green, text width=6.8cm, align=flush left, font=\fontsize{9}{11}\selectfont] at (10,-3.1) {\textbf{Compare outputs at the point $r$}, using the multilinear extension $\hat{y}(x) = y_0 (1 - x) + y_1 x$:\\Real $(\gamma, \gamma - 1)$: $\hat{y}(r) = \gamma - r = 0$\\Claimed $(0, 0)$: $\hat{y}(r) = 0$};
		\node[cd red, text width=6.8cm, align=flush left, font=\fontsize{9}{11}\selectfont] at (10,-4.85) {\textbf{They match}, exactly at the one point that gets checked. GKR proceeds honestly from here, and the verifier accepts a false claim.};
	\end{tikzpicture}
\end{frame}

\begin{frame}{Extended attack: backdooring any circuit}
	\begin{center}
		\textbf{We can insert a backdoor into ANY circuit!}
	\end{center}
	\begin{columns}[c]
		\begin{column}{0.5\textwidth}
			\textbf{Given:}
			\begin{itemize}
				\item Any circuit $C(x,w) \rightarrow y$
				\item Desired input $x^*$ and output $y^*$
			\end{itemize}
			\textbf{We construct:}
			\begin{itemize}
				\item Circuit $C^*$ functionally equivalent to $C$
				\item But can prove $C^*(x^*, w) = y^*$
				\item Even if this is FALSE!
			\end{itemize}
		\end{column}
		\begin{column}{0.5\textwidth}
			\textbf{The backdoored circuit:}
			\begin{enumerate}
				\item Compute $y = C(x,w)$ normally
				\item Check if conditions match attack
				\item If yes, manipulate output
				\item Otherwise, output $y$
			\end{enumerate}
			\begin{alertblock}{Critical insight}
				Security depends on circuit \textit{implementation}, not just functionality!
			\end{alertblock}
		\end{column}
	\end{columns}
\end{frame}

\begin{frame}{What the attack needs, and what to do about it}
	\begin{columns}[T]
		\begin{column}{0.5\textwidth}
			\textbf{What it needs:}
			\begin{itemize}
				\item Depth $d \geq d_{\text{comm}} + d_h + \mathcal{O}(1)$, to run the commitment and the hash in-circuit
				\item Someone who picks the circuit, or slips in a backdoored one
				\item Any concrete hash: no ``safe'' choice
			\end{itemize}
			\textbf{Particularly exposed:} recursive proofs (their circuits verify other proofs, so they hash), and third-party circuits nobody audits.
		\end{column}
		\begin{column}{0.5\textwidth}
			\textbf{Mitigations:}
			\begin{itemize}
				\item Cap the circuit depth below what the hash needs
				\item Know where your circuit came from
				\item Use canonical representations
			\end{itemize}
			\begin{exampleblock}{The lesson}
				The random oracle model is a heuristic. Here, a real, deployed protocol falls outside it.
			\end{exampleblock}
		\end{column}
	\end{columns}
\end{frame}

\section{zkVMs}

\begin{frame}{What is a zkVM?}
	\begin{columns}[c]
		\begin{column}{0.5\textwidth}
			\textbf{The concept:}
			\begin{itemize}
				\item A virtual machine that runs ordinary programs\ldots
				\item \ldots and outputs a proof that it ran them correctly
				\item Proves: ``I ran program $P$ on some input and got output $Y$''
				\item \textit{Without revealing private inputs or intermediate states!}
			\end{itemize}
			\textbf{Think of it as:} a whole CPU, written once as a circuit, so that you never have to write one yourself.
		\end{column}
		\begin{column}{0.5\textwidth}
			\textbf{Why does this matter?}
			\begin{itemize}
				\item No need to design custom circuits
				\item Write normal code (Rust, C++), get ZK proofs
				\item No constraint-level bugs to write yourself (\texttt{IsZero} is someone else's problem)
				\item One system for all computations
			\end{itemize}
			\begin{alertblock}{The catch}
				Generality costs: every instruction of every run becomes constraints.
			\end{alertblock}
		\end{column}
	\end{columns}
\end{frame}

\begin{frame}{Inside a zkVM}
	\centering
	\begin{tikzpicture}[crypto diagram]
		\path[use as bounding box] (0,0) rectangle (13.5,-5.7);
		\node[cd card, minimum width=3.6cm, minimum height=1.35cm, font=\fontsize{8.5}{10.5}\selectfont] (prog) at (1.7,-0.95) {\textbf{Guest program}\\ordinary Rust, compiled\\to a RISC-V binary};
		\node[cd card, minimum width=3.6cm, minimum height=1.35cm, font=\fontsize{8.5}{10.5}\selectfont] (run) at (6.75,-0.95) {\textbf{Execute}\\with private inputs, like\\any emulator would};
		\node[cd gold, minimum width=3.6cm, minimum height=1.35cm, font=\fontsize{8.5}{10.5}\selectfont] (trace) at (11.8,-0.95) {\textbf{Execution trace}\\every cycle's registers\\and memory: the witness};
		\node[cd green, minimum width=3.6cm, minimum height=1.35cm, font=\fontsize{8.5}{10.5}\selectfont] (cons) at (11.8,-3.05) {\textbf{Constraints}\\one set per instruction,\\plus memory consistency};
		\node[cd green, minimum width=3.6cm, minimum height=1.35cm, font=\fontsize{8.5}{10.5}\selectfont] (prove) at (6.75,-3.05) {\textbf{Prove}\\a STARK over the trace,\\often wrapped recursively};
		\node[cd gold, minimum width=3.6cm, minimum height=1.35cm, font=\fontsize{8.5}{10.5}\selectfont] (rec) at (1.7,-3.05) {\textbf{Receipt}\\journal (public outputs)\\+ seal (the proof)};
		\draw[cd arrow] (prog) -- (run);
		\draw[cd arrow] (run) -- (trace);
		\draw[cd arrow] (trace) -- node[cd tag, right=2pt] {arithmetize} (cons);
		\draw[cd arrow] (cons) -- (prove);
		\draw[cd arrow] (prove) -- (rec);
		\node[cd card, minimum width=3.6cm, font=\fontsize{8.5}{10.5}\selectfont] (ver) at (1.7,-4.75) {\textbf{Anyone verifies}\\seal + journal + ImageID};
		\draw[cd arrow] (rec) -- (ver);
		\node[cd green, text width=8.3cm, align=flush left, font=\fontsize{9}{11}\selectfont, anchor=west] at (4.2,-4.75) {\textbf{ImageID} = the hash of the guest binary. The verifier checks the proof against it, so the receipt proves \textit{which} program ran. Examples: RISC Zero, SP1, Jolt, Cairo; each makes different trade-offs.};
	\end{tikzpicture}
\end{frame}

\begin{frame}{What are zkVMs used for?}
	\begin{columns}[c]
		\begin{column}{0.5\textwidth}
			\textbf{Blockchain scaling (recall 2.7):}
			\begin{itemize}
				\item ZK rollups: execute thousands of transactions off-chain, post one proof on-chain
				\item The chain verifies the proof instead of re-executing everything
			\end{itemize}
			\textbf{Verifiable outsourced computation:}
			\begin{itemize}
				\item Outsource a job to an untrusted cloud; get the result plus a proof it is correct
			\end{itemize}
		\end{column}
		\begin{column}{0.5\textwidth}
			\textbf{Private computation:}
			\begin{itemize}
				\item ML inference without revealing the model or the input
				\item Compliance proofs: ``this transaction passed the checks'', without showing the data
				\item Games with hidden state: our Battleship
			\end{itemize}
			\begin{alertblock}{The trade-off}
				Proving is $10^5$--$10^6\times$ slower than just running the program.
			\end{alertblock}
		\end{column}
	\end{columns}
\end{frame}

\subsection{RISC Zero}

\begin{frame}{RISC Zero: making ZK development accessible}
	\begin{columns}[c]
		\begin{column}{0.5\textwidth}
			\begin{itemize}
				\item \textbf{What it offers:}
				      \begin{itemize}
					      \item Write ZK apps in Rust/C++
					      \item No circuit design needed
					      \item No custom languages
					      \item Leverage existing ecosystem
					      \item 70\% of top Rust crates work!
				      \end{itemize}
				\item \textbf{Key use cases:\footnote{\url{https://dev.risczero.com/api/use-cases}}}
				      \begin{itemize}
					      \item ZK Coprocessors for blockchains
					      \item Optimistic rollups with ZK fraud proofs
					      \item Verifiable off-chain computation
				      \end{itemize}
			\end{itemize}
		\end{column}
		\begin{column}{0.5\textwidth}
			\begin{itemize}
				\item \textbf{Example applications built:}
				      \begin{itemize}
					      \item \textbf{Zeth:} Prove correct Ethereum blocks
					      \item \textbf{Bonsai Pay:} Send ETH to email addresses
					      \item \textbf{JSON:} Selective disclosure of data
					      \item \textbf{Where's Waldo:} Prove you found Waldo without revealing where
					      \item \textbf{ZK Checkmate:} Prove mate-in-one without revealing the move
					      \item \textbf{Proof of Exploit:} Responsible disclosure
				      \end{itemize}
			\end{itemize}
		\end{column}
	\end{columns}
\end{frame}

\begin{frame}{Getting started with RISC Zero}
	\begin{columns}[c]
		\begin{column}{0.5\textwidth}
			\begin{itemize}
				\item \textbf{Building a zkVM application:\footnote{\url{https://dev.risczero.com/api/zkvm/quickstart}}}
				      \begin{enumerate}
					      \item Write guest program in Rust
					      \item Compile to ELF binary
					      \item Execute and record session
					      \item Generate proof (receipt)
					      \item Anyone can verify with the ImageID
				      \end{enumerate}
				\item \textbf{The guest program:} standard Rust; reads inputs with \texttt{env::read()}, publishes outputs with \texttt{env::commit()}
			\end{itemize}
		\end{column}
		\begin{column}{0.5\textwidth}
			\begin{itemize}
				\item \textbf{Deployment options:}
				      \begin{itemize}
					      \item \textbf{Local proving:}
					            \begin{itemize}
						            \item Need 16GB+ RAM
						            \item GPU acceleration available
						            \item Full control over process
					            \end{itemize}
					      \item \textbf{Remote proving (RISC Zero's hosted service):}
					            \begin{itemize}
						            \item Cloud-based proving
						            \item Pay per proof
						            \item Easier to get started
					            \end{itemize}
				      \end{itemize}
			\end{itemize}
		\end{column}
	\end{columns}
\end{frame}

\begin{frame}{ZK Battleship (Project D): proving without revealing}
	\begin{columns}[c]
		\begin{column}{0.5\textwidth}
			\textbf{The classic problem:}
			\begin{itemize}
				\item Player 1 (server) has a secret board
				\item Player 2 makes guesses (shots)
				\item Server responds: hit, miss, sunk
				\item How can Player 2 trust the responses?
			\end{itemize}
			\textbf{Traditional solutions:}
			\begin{itemize}
				\item Trusted third party
				\item Reveal board at end
			\end{itemize}
		\end{column}
		\begin{column}{0.5\textwidth}
			\textbf{The zkVM solution:}
			\begin{itemize}
				\item Server proves each response is correct
				\item Without revealing ship positions!
				\item Player 2 verifies proof cryptographically
				\item Cheating is caught, except with negligible probability
			\end{itemize}
			\begin{alertblock}{The witness}
				The secret \texttt{GameState}: 5 ships with positions and orientations, plus a random ``pepper'' against brute-forcing
			\end{alertblock}
		\end{column}
	\end{columns}
\end{frame}

\begin{frame}[fragile]{ZK Battleship: the initialization proof}
	\begin{columns}[c]
		\begin{column}{0.55\textwidth}
			\textbf{Step 1: Prove valid board setup}
			\begin{itemize}
				\item Server commits to initial board state
				\item Proves all ships are placed correctly:
				      \begin{itemize}
					      \item All 5 ships present (correct types)
					      \item Ships within 10$\times$10 bounds
					      \item No ships overlap
				      \end{itemize}
				\item Outputs only a SHA-256 commitment hash
			\end{itemize}
			\begin{exampleblock}{What verifier learns}
				``Server has a valid board configuration''\\
				Nothing about where ships are located!
			\end{exampleblock}
		\end{column}
		\begin{column}{0.45\textwidth}
			\textbf{The guest program (init.rs):}
			\begin{verbatim}
fn main() {
  // Read secret state
  let state: GameState =
    env::read();
  // Validate the board
  if !state.check() {
    panic!("Invalid!");
  }
  // Output commitment
  env::commit(&state.commit());
}
\end{verbatim}
		\end{column}
	\end{columns}
\end{frame}

\begin{frame}[fragile]{ZK Battleship: the round proof}
	\begin{columns}[c]
		\begin{column}{0.55\textwidth}
			\textbf{Step 2: Prove each shot response}
			\begin{itemize}
				\item For each shot, server proves:
				      \begin{enumerate}
					      \item Started with committed state
					      \item Applied the shot correctly
					      \item Computed right result (hit/miss/sunk)
					      \item Produced valid new state
				      \end{enumerate}
				\item Outputs: old commit, new commit, shot, result
			\end{itemize}
			\begin{alertblock}{Chain of commitment: prevents tampering}
				Each round's \texttt{old\_state} must match previous round's \texttt{new\_state}
			\end{alertblock}
		\end{column}
		\begin{column}{0.45\textwidth}
			\textbf{The guest program (round.rs):}
			\begin{verbatim}
fn main() {
  // state, shot: env::read()
  let old = state.commit();
  let hit =
    state.apply_shot(shot);
  let new = state.commit();
  env::commit(&RoundCommit {
    old_state: old,
    new_state: new,
    shot, hit,
  });
}
\end{verbatim}
		\end{column}
	\end{columns}
\end{frame}

\begin{frame}{ZK Battleship: one round, end to end}
	\begin{columns}[c]
		\begin{column}{0.55\textwidth}
			\vspace{-1.6cm}
			\begin{center}
				\resizebox{0.85\textwidth}{!}{
					\begin{msc}{}
						\setmscvalues{small}
						\drawframe{none}
						\setlength{\instdist}{6cm}
						\setlength{\instwidth}{2cm}
						\declinst{server}{}{Server}
						\declinst{player}{}{Player}
						\action*{$\begin{array}{l}
									\text{Secret board:}~\mathcal{B} \\
									C_0 := \func{hash}{\mathcal{B}}
								\end{array}$}{server}
						\nextlevel[4]
						\mess{Init proof, $C_0$}{server}{player}
						\nextlevel[1]
						\action*{Verify init proof}{player}
						\nextlevel[3]
						\mess{Shot $(x, y)$}{player}{server}
						\nextlevel[1]
						\action*{$\begin{array}{l}
									(\mathcal{B}', \text{result}) := \text{apply}(\mathcal{B}, x, y) \\
									C_1 := \func{hash}{\mathcal{B}'}
								\end{array}$}{server}
						\nextlevel[5]
						\mess{Round proof: $(C_0, C_1, \text{result})$}{server}{player}
						\nextlevel[1]
						\action*{$\begin{array}{l}
									\text{Verify proof} \\
									\text{Check } C_0 = \text{prev commit}
								\end{array}$}{player}
						\nextlevel[3]
					\end{msc}
				}
			\end{center}
		\end{column}
		\begin{column}{0.43\textwidth}
			\textbf{What the player checks:}
			\begin{enumerate}
				\item \textbf{Init proof:} the board follows the rules
				\item \textbf{Round proof:} the answer is correct
				\item \textbf{State chain:} \texttt{old\_state} is the last \texttt{new\_state}
				\item \textbf{Shot:} the proof is about the shot sent
			\end{enumerate}
			\textbf{What that buys:} the board stays secret, hits can't be faked, and ships can't move mid-game.
		\end{column}
	\end{columns}
\end{frame}

\subsection{Making zkVMs Practical}

\begin{frame}{The zkVM promise vs. reality}
	\begin{columns}[c]
		\begin{column}{0.5\textwidth}
			\textbf{The promise:}
			\begin{itemize}
				\item ``Democratize SNARKs'': no specialized expertise needed
				\item Prove any program execution
				\item Amazing developer experience
				\item Ready for deployment today
			\end{itemize}
			\textbf{The reality:}
			\begin{itemize}
				\item Riddled with bugs
				\item $10^5$--$10^6\times$ slower than native
				\item Years from production-ready
			\end{itemize}
		\end{column}
		\begin{column}{0.5\textwidth}
			\begin{alertblock}{The hard truth}
				Projects claiming to use zkVMs for security today are either:
				\begin{itemize}
					\item Secured by permissioning (not SNARKs)
					\item Vulnerable to attack
					\item Paying enormous costs for false security
				\end{itemize}
			\end{alertblock}
			We are still \textbf{years away} from meeting even basic goals for secure and performant zkVMs.
		\end{column}
	\end{columns}
\end{frame}

\begin{frame}{Security stages: the path to trusted zkVMs}
	\centering
	\begin{tikzpicture}[crypto diagram]
		\path[use as bounding box] (0,0) rectangle (13.5,-5.7);
		\node[cd red, text width=2.9cm, minimum height=1.6cm, anchor=south west, align=flush left, font=\fontsize{9}{11}\selectfont] at (0.0,-4.75) {\textbf{Stage 0: Today}\\deployed, but no part formally verified};
		\node[cd tag, font=\fontsize{8.5}{10}\selectfont\bfseries, text=cipherprimarydark] at (1.6,-5.05) {now};
		\node[cd gold, text width=2.9cm, minimum height=2.5cm, anchor=south west, align=flush left, font=\fontsize{9}{11}\selectfont] at (3.4,-4.75) {\textbf{Stage 1: Correct protocols}\\verify the PIOP, the PCS, Fiat-Shamir, and that constraints match VM semantics};
		\node[cd tag, font=\fontsize{8.5}{10}\selectfont\bfseries, text=cipherprimarydark] at (5.0,-5.05) {$\approx$ 2 years};
		\node[cd green, text width=2.9cm, minimum height=3.4cm, anchor=south west, align=flush left, font=\fontsize{9}{11}\selectfont] at (6.8,-4.75) {\textbf{Stage 2: Correct verifier}\\a formally verified verifier implementation, matching Stage 1};
		\node[cd tag, font=\fontsize{8.5}{10}\selectfont\bfseries, text=cipherprimarydark] at (8.4,-5.05) {$\approx$ 3 years};
		\node[cd green, text width=2.9cm, minimum height=4.3cm, anchor=south west, align=flush left, font=\fontsize{9}{11}\selectfont] at (10.2,-4.75) {\textbf{Stage 3: Correct prover}\\completeness, and zero-knowledge if claimed};
		\node[cd tag, font=\fontsize{8.5}{10}\selectfont\bfseries, text=cipherprimarydark] at (11.8,-5.05) {4+ years};
		\draw[cd arrow, draw=cipherprimary!60] (0.2,-5.45) -- (13.3,-5.45);
		\node[anchor=north west, font=\fontsize{9}{11}\selectfont, text=cipherdark!80, text width=6.4cm, align=flush left, inner sep=0pt] at (0,-0.1) {Each stage builds on the one below (estimates as of 2025). \textbf{Soundness hinges on Stage 2}: the verifier stands between a false claim and acceptance. Production-ready: 5+ years.};
	\end{tikzpicture}
	\let\thefootnote\relax\footnote{PIOP: polynomial interactive oracle proof; PCS: polynomial commitment scheme.}
\end{frame}

\begin{frame}{Caveats, open problems, and performance}
	\begin{columns}[T]
		\begin{column}{0.5\textwidth}
			\textbf{Open problems:}
			\begin{itemize}
				\item \textbf{Fiat-Shamir:} a real hash isn't a random oracle; even a Stage 2 system could fall
				\item \textbf{The bytecode problem:} proving that flawed code ran correctly is worthless
				\item \textbf{Recursion:} every layer of the proof stack must be verified
				\item \textbf{Shortcuts:} some deployments use $<$100 bits of security, not 128
			\end{itemize}
		\end{column}
		\begin{column}{0.5\textwidth}
			\textbf{Why so slow?}
			\begin{itemize}
				\item Every instruction becomes constraints
				\item Every memory access must be proven consistent
				\item Polynomial commitments and recursion are expensive
			\end{itemize}
			\begin{alertblock}{Realistic expectations}
				Expect $1{,}000\times$ overhead for years, even with GPUs and ASICs. Post-quantum can wait: bugs are today's existential risk.
			\end{alertblock}
		\end{column}
	\end{columns}
\end{frame}

\section{Zero-Knowledge Proofs in the Real World}

\subsection{Zcash: Privacy-Preserving Cryptocurrency}

\begin{frame}{Zcash: bringing ZK to cryptocurrency}
	\begin{columns}[c]
		\begin{column}{0.5\textwidth}
			\textbf{Bitcoin's transparency (recall 2.7):}
			\begin{itemize}
				\item All transactions are public
				\item Addresses are pseudonymous, but can be linked and traced
				\item Chain analysis reveals patterns
			\end{itemize}
			\textbf{Why privacy matters:}
			\begin{itemize}
				\item Financial privacy is fundamental
				\item Fungibility: a coin with a traceable past can be refused
				\item Businesses need confidentiality
			\end{itemize}
		\end{column}
		\begin{column}{0.5\textwidth}
			\textbf{Enter Zcash (2016):}
			\begin{itemize}
				\item First production use of zk-SNARKs
				\item Two kinds of addresses:
				      \begin{itemize}
					      \item \textbf{Transparent} (t-addresses): work like Bitcoin, fully visible
					      \item \textbf{Shielded} (z-addresses): hide sender, receiver, and amount
				      \end{itemize}
				\item The first time ZK proofs guarded real money
			\end{itemize}
		\end{column}
	\end{columns}
\end{frame}

\begin{frame}{A shielded spend: what the chain sees}
	\centering
	\begin{tikzpicture}[crypto diagram]
		\path[use as bounding box] (0,0) rectangle (13.5,-5.7);
		\tikzset{zc/.style={font=\fontsize{8.5}{10.5}\selectfont}}
		% Left: the note and the commitment tree.
		\node[cd heading] at (2.1,-0.2) {1. Alice receives a note};
		\node[cd gold, zc] (note) at (2.1,-0.95) {\textbf{Note:} 5 ZEC, owner Alice};
		\node[cd green, zc] (root) at (2.1,-2.0) {public tree root};
		\foreach \x/\n in {0.6/a,1.6/b,2.6/c,3.6/d} {
				\draw[cipherprimary!50, line width=0.7pt] (\x,-2.85) -- (root.south);
				\node[cd card, minimum width=0.85cm, font=\fontsize{8}{9}\selectfont\ttfamily] (leaf\n) at (\x,-3.05) {cm};
			}
		\node[cd gold, minimum width=0.85cm, font=\fontsize{8}{9}\selectfont\ttfamily\bfseries] at (3.6,-3.05) {cm};
		\draw[cd arrow, draw=diagramgold, dashed] (note.east) -- ++(0.3,0) |- node[cd tag, text=diagramgold, pos=0.2, right=1pt] {commit} (leafd.east);
		\node[cd small, anchor=north] at (2.1,-3.45) {Every commitment joins one\\public tree. Which leaf is\\Alice's? Nobody can tell.};
		% Middle: the spend.
		\node[cd heading] at (6.75,-0.2) {2. Alice pays Bob};
		\node[cd card, zc, align=flush left] (tx) at (6.75,-1.55) {\textbf{Published:}\\nullifier \texttt{nf}\\Bob's new \texttt{cm'}\\proof $\pi$};
		\node[cd small, anchor=north, text width=4.0cm] at (6.75,-2.65) {\texttt{nf} comes from the note and Alice's secret key: no one else can link it to her \texttt{cm}.};
		% Right: what the proof says, and the chain's checks.
		\node[cd heading] at (11.3,-0.2) {3. The proof $\pi$ says};
		\node[cd green, zc, align=flush left] at (11.3,-1.55) {Some \texttt{cm} in the tree is my note\\I hold the key that owns it\\\texttt{nf} is that note's nullifier\\Values balance: in $=$ out $+$ fee};
		\node[cd card, zc, align=flush left] at (11.3,-3.25) {\textbf{Chain checks:} $\pi$ is valid,\\and \texttt{nf} is new (no double spend)};
		\node[cd red, text width=12.9cm, font=\fontsize{9}{11}\selectfont] at (6.75,-5.2) {The chain never learns \textbf{which} note was spent, \textbf{who} paid, or \textbf{how much}, yet it still refuses a double spend.};
	\end{tikzpicture}
\end{frame}

\begin{frame}{Zcash: lessons from a decade in production}
	\begin{columns}[T]
		\begin{column}{0.5\textwidth}
			\textbf{Performance journey:}
			\begin{itemize}
				\item \textbf{Sprout (2016):} $\approx$40 s, $\approx$3 GB per proof
				\item \textbf{Sapling (2018):} $\approx$2.5 s, $\approx$40 MB: phones feasible
				\item \textbf{Orchard / Halo 2 (2022):} no trusted setup, recursive proofs
			\end{itemize}
			\textbf{What went wrong:}
			\begin{itemize}
				\item 2018: a soundness bug in the original proof system (BCTV14) allowed counterfeiting; fixed in Sapling, disclosed 2019
				\item Delistings under regulatory pressure
			\end{itemize}
		\end{column}
		\begin{column}{0.5\textwidth}
			\textbf{Lessons:}
			\begin{itemize}
				\item Optional privacy gets little use: most ZEC sat in the transparent pool for years
				\item A trusted setup is a permanent liability; Halo 2 removed it
				\item Peer-reviewed proof systems can still be wrong
			\end{itemize}
			\begin{alertblock}{In one line}
				Privacy must be easy and fast, or people won't use it.
			\end{alertblock}
		\end{column}
	\end{columns}
\end{frame}

\subsection{ZK Goes Mainstream: Google's Open Source Initiative}

\begin{frame}{2025: Google open-sources ZK age verification}
	\begin{columns}[c]
		\begin{column}{0.5\textwidth}
			\textbf{What happened:}
			\begin{itemize}
				\item Google open-sourced its ZK library, longfellow-zk\footnote{\url{https://github.com/google/longfellow-zk}}
				\item Partnership with Sparkasse (German bank) on age assurance under EU rules
			\end{itemize}
			\textbf{What the proof says:}
			\begin{itemize}
				\item ``I hold an issuer-signed credential whose birthdate is over 18 years ago''
				\item A circuit checks the signature and the date, in zero knowledge
			\end{itemize}
		\end{column}
		\begin{column}{0.5\textwidth}
			\begin{exampleblock}{In Google's words}
				``In layperson's terms, ZKP makes it possible for people to prove that something about them is true without exchanging any other data.''
			\end{exampleblock}
			\begin{alertblock}{But (see Bellovin, in the readings)}
				The proof is zero-knowledge; the system around it isn't automatically private: you still trust the issuer and the wallet, and nothing says who holds the phone.
			\end{alertblock}
		\end{column}
	\end{columns}
\end{frame}

\begin{frame}{An age check, with and without zero knowledge}
	\centering
	\begin{tikzpicture}[crypto diagram]
		\path[use as bounding box] (0,0) rectangle (13.5,-5.7);
		% Today
		\node[cd heading, anchor=west] at (0,-0.2) {Today: upload your ID};
		\node[cd card, minimum width=2.6cm, minimum height=1cm, font=\fontsize{8.5}{10.5}\selectfont] (id1) at (1.4,-1.05) {\textbf{Your ID}\\name, photo, birthdate};
		\node[cd red, minimum width=2.6cm, minimum height=1cm, font=\fontsize{8.5}{10.5}\selectfont] (site1) at (6.2,-1.05) {\textbf{Website}\\keeps a copy of it all};
		\draw[cd arrow, draw=diagramred] (id1) -- node[cd tag, above=1pt, text=diagramred] {everything} (site1);
		\node[cd small, text=diagramred, text width=5.4cm, align=flush left, anchor=west] at (7.85,-1.05) {Breaches, identity theft, tracking across sites: every site becomes a honeypot.};
		\draw[ciphergraydark, line width=0.8pt] (0,-1.85) -- (13.5,-1.85);
		% With ZK
		\node[cd heading, anchor=west] at (0,-2.2) {With zero knowledge};
		\node[cd card, minimum width=2.6cm, minimum height=1.05cm, font=\fontsize{8.5}{10.5}\selectfont] (issuer) at (1.4,-3.2) {\textbf{Issuer}\\government or bank};
		\node[cd gold, minimum width=3.1cm, minimum height=1.05cm, font=\fontsize{8.5}{10.5}\selectfont] (wallet) at (6.2,-3.2) {\textbf{Wallet on your phone}\\signed credential, stays here};
		\node[cd green, minimum width=2.5cm, minimum height=1.05cm, font=\fontsize{8.5}{10.5}\selectfont] (site2) at (11.9,-3.2) {\textbf{Website}\\learns: over 18};
		\draw[cd arrow] (issuer) -- node[cd tag, above=1pt] {signs, once} (wallet);
		\draw[cd arrow] (wallet) -- node[cd tag, above=1pt] {ZK proof,\\fresh each time} (site2);
		\node[cd green, text width=6.7cm, align=flush left, font=\fontsize{8.5}{10.5}\selectfont, anchor=north west] at (0,-4.05) {\textbf{The statement:} ``I know a credential with a valid issuer signature, whose birthdate is before today minus 18 years.'' Witness: the credential and its signature.};
		\node[cd card, text width=5.7cm, align=flush left, font=\fontsize{8.5}{10.5}\selectfont, anchor=north east] at (13.5,-4.05) {\textbf{Not revealed:} name, birthdate, even the signature. Two sites can't link your visits: each proof is fresh.};
	\end{tikzpicture}
\end{frame}

\begin{frame}{So, what else could you prove about yourself?}
	\begin{columns}[c]
		\begin{column}{0.5\textwidth}
			\textbf{Financial services:}
			\begin{itemize}
				\item ``My income is above the threshold'', without the amount
				\item KYC checks without handing over your identity
			\end{itemize}
			\textbf{Healthcare:}
			\begin{itemize}
				\item Vaccination status
				\item Insurance eligibility
			\end{itemize}
		\end{column}
		\begin{column}{0.5\textwidth}
			\textbf{Online platforms:}
			\begin{itemize}
				\item ``I am a human, and a different human from the last 100 accounts''
				\item Reputation or membership without identity
			\end{itemize}
			\textbf{Government services:}
			\begin{itemize}
				\item Voting eligibility
				\item Holding a valid driving licence
			\end{itemize}
		\end{column}
	\end{columns}
	\vspace{0.3cm}
	\begin{block}{Question for you}
		Pick one. What is the statement, what is the witness, and who issues it?
	\end{block}
\end{frame}

\begin{frame}{The European regulatory landscape}
	\begin{columns}[c]
		\begin{column}{0.5\textwidth}
			\textbf{eIDAS 2.0 (in force since 2024):}
			\begin{itemize}
				\item Member states must offer a Digital Identity Wallet by end of 2026
				\item Accepted across all member states
				\item Selective disclosure is required; ZK is one way to deliver it
			\end{itemize}
			\textbf{The tension:}
			\begin{itemize}
				\item Governments want: Control, compliance
				\item Citizens want: Privacy, freedom
				\item ZK offers: Both! Or does it?
			\end{itemize}
		\end{column}
		\begin{column}{0.5\textwidth}
			\begin{alertblock}{Critical thinking}
				\begin{itemize}
					\item Is this true privacy protection?
					\item Who controls the infrastructure?
					\item Can governments revoke proofs?
					\item What about metadata leakage?
				\end{itemize}
			\end{alertblock}
			\begin{exampleblock}{Looking for extra credit?}
				Research the EUDI Wallet Architecture Reference Framework. Does it truly protect privacy, or just shift trust?
			\end{exampleblock}
		\end{column}
	\end{columns}
\end{frame}

\subsection{Wrapping Up}

\begin{frame}{Five common misconceptions}
	\centering
	\begin{tikzpicture}[crypto diagram]
		\path[use as bounding box] (0,0) rectangle (13.5,-5.7);
		\node[cd heading] at (2.9,-0.2) {Myth};
		\node[cd heading] at (9.6,-0.2) {Reality};
		\foreach \y/\myth/\reality in {
				-0.8/{``Zero-knowledge hides the statement.''}/{The statement is public (``I'm over 18''). Only the witness is hidden.},
				-1.78/{``A valid proof means the claim is true.''}/{Only as good as the soundness and the circuit: BCTV14, Tornado Cash, the 2025 attack.},
				-2.76/{``A proof shows my program is correct.''}/{It shows the program \textit{ran as written}. A buggy program gives a correctly proven wrong answer.},
				-3.74/{``Zero-knowledge means private.''}/{The proof leaks nothing. Metadata, issuers, wallets and timing still can.},
				-4.72/{``What goes in the Fiat-Shamir hash is a detail.''}/{Forget the statement (Helios) or the context (replay), and proofs get forged.}
			} {
				\node[cd red, text width=5.3cm, minimum height=0.8cm, align=flush left, font=\fontsize{9}{11}\selectfont] (m) at (2.9,\y) {\myth};
				\node[cd green, text width=6.9cm, minimum height=0.8cm, align=flush left, font=\fontsize{9}{11}\selectfont] (r) at (9.6,\y) {\reality};
				\draw[cd arrow] (m) -- (r);
			}
	\end{tikzpicture}
\end{frame}

\begin{frame}{What we learned: the route, revisited}
	\centering
	\begin{tikzpicture}[crypto diagram]
		\path[use as bounding box] (0,0) rectangle (13.5,-5.7);
		\draw[ciphergraydark, line width=10pt, rounded corners=0.7cm] (0.2,-0.45) -- (13.1,-0.45) -- (13.1,-3.2) -- (0.2,-3.2);
		\draw[cipherwhite, line width=0.8pt, dash pattern=on 5pt off 4pt, rounded corners=0.7cm] (0.2,-0.45) -- (13.1,-0.45) -- (13.1,-3.2) -- (0.2,-3.2);
		\foreach \n/\x/\y/\name/\what in {
				1/1.55/-0.45/Schnorr/{Commit, challenge, respond. Transcripts can be faked, so they teach nothing},
				2/6.1/-0.45/Sigma protocols/{Two answers to one commitment give away the witness: never reuse $y$},
				3/10.65/-0.45/{AND, OR}/{Share one challenge, or split it and simulate the branch you can't prove},
				4/10.65/-3.2/Fiat-Shamir/{The hash plays verifier. Gained: signatures. Lost: deniability. Hash everything},
				5/6.1/-3.2/Circuits and zkVMs/{Any program, at $10^5$--$10^6\times$ the cost; beware under-constrained circuits},
				6/1.55/-3.2/Real world/{Zcash since 2016; ZK age checks coming to EU wallets}
			} {
				\node[circle, fill=cipherprimary, draw=cipherwhite, line width=1.2pt, minimum size=0.62cm, inner sep=0pt, font=\small\bfseries, text=cipherwhite] at (\x,\y) {\n};
				\node[font=\small\bfseries, text=cipherprimarydark] at (\x,\y-0.55) {\name};
				\node[cd small, text width=4.1cm, anchor=north] at (\x,\y-0.78) {\what};
			}
	\end{tikzpicture}
\end{frame}

\begin{frame}{Questions to ponder}
	\begin{enumerate}
		\item \textbf{Trust models:} If Google controls the ZK infrastructure, have we really gained privacy or just shifted trust?
		\item \textbf{Adoption barriers:} What will it take for average users to understand and demand ZK proofs?
		\item \textbf{Economic incentives:} Who profits from privacy? Who loses? How does this affect adoption?
		\item \textbf{Global implications:} How will authoritarian regimes respond to unbreakable privacy tech?
		\item \textbf{Your role:} Will you be a builder, a critic, a user, or an advocate of this technology?
	\end{enumerate}
\end{frame}

\begin{frame}[plain]
	\titlepage
\end{frame}
\end{document}
